
\documentclass{article} % For LaTeX2e
\usepackage{iclr2027_conference,times}
\usepackage{enumitem}
\usepackage{graphicx}
\usepackage{booktabs}
\usepackage{tabularx}

% Optional math commands from https://github.com/goodfeli/dlbook_notation.
\input{math_commands.tex}

\usepackage{hyperref}
\usepackage{url}


\title{Truncate to Accelerate: 
\\From Forecast Recoverability to Efficient \\Time-Series Foundation Models} 
% Truncate to Accelerate: Depth, Geometry, and Redundancy in Time-Series Foundation Models 


\newcommand{\fix}{\marginpar{FIX}}
\newcommand{\new}{\marginpar{NEW}}
%\iclrfinalcopy % Uncomment for camera-ready version, but NOT for submission.

\begin{document}


\maketitle
\begin{abstract}
In large language models, intermediate representations have been used both to study how task-relevant information evolves with depth and to accelerate inference. In time-series foundation models (TSFMs), however, these questions remain largely unexplored. We address this gap through layerwise probing, representation geometry, and lightweight adaptation across three state-of-the-art TSFMs: Chronos-2, TimesFM-3, and TiRex, evaluated on a collection of datasets with diverse sampling frequencies. First, layerwise forecasting readouts reveal a forecasting tunnel: across nearly all model-dataset pairs, intermediate depths already recover performance within 5\% of the final-block readout. Next, we find that functional saturation coexists with continued geometric evolution. Finally, intermediate states can contain recoverable forecast information, but feeding them directly into the model's native forecasting pathway degrades performance relative to the full model. Lightweight affine alignment substantially reduces this gap and exposes a measured performance-compute frontier: aligned Chronos-2 remains within a median 3.1\% of the full model after retaining only 4 of its 12 blocks, while running $2.43\times$ faster, with the achievable tradeoff varying substantially across architectures. Moreover, cuts inside the forecasting tunnel are typically better tolerated at matched depth, linking early recoverability to practical truncation. Together, these results suggest that much of the forecast-relevant information becomes recoverable before the final block and offer a path toward both understanding and accelerating time-series foundation models.
\end{abstract}


\section{Introduction}
\label{sec:introduction}

Time-series foundation models (TSFMs) have recently achieved strong zero-shot forecasting performance across a broad range of domains
\citep{ansari2024chronoslearninglanguagetime,
das2024decoderonlyfoundationmodeltimeseries,
woo2024unifiedtraininguniversaltime,
ansari2025chronos2univariateuniversalforecasting,
auer2025tirexzeroshotforecastinglong}.
This capability increasingly relies on deep and computationally expensive backbones. At inference time, these models typically execute their full sequence of representation-transforming stages before producing a forecast. It is not clear, however, that every one of these stages still improves forecasting performance; identifying where performance saturates could therefore provide a direct route to faster inference.

In large language models, substantial depth can be removed with limited degradation \citep{fan2019reducingtransformerdepthdemand, gromov2025unreasonableineffectivenessdeeperlayers, men2024shortgptlayerslargelanguage}. For TSFMs, \citet{wiliński2025exploringrepresentationsinterventionstime} locate removable blocks through similarity between intermediate representations. Representational similarity, however, need not coincide with forecasting redundancy. We therefore ask the functional question directly: \emph{at what depth does near-final forecasting performance become recoverable, and can that early recoverability be converted into actual truncation?}

We study this in three state-of-the-art TSFMs with distinct architectures: the encoder-only transformer Chronos-2 \citep{ansari2025chronos2univariateuniversalforecasting}, the decoder-only transformer TimesFM-3 \citep{das2024decoderonlyfoundationmodeltimeseries}, and the xLSTM-based TiRex \citep{auer2025tirexzeroshotforecastinglong}, across 14 datasets spanning five sampling frequencies. The analysis proceeds in four steps. First, we introduce the \emph{sustained forecasting tunnel}, a task-level criterion for locating the depth at which layerwise forecasting readouts reach near-final performance. We then compare this functional notion of depth with representation geometry, using debiased linear Centered Kernel Alignment (CKA) \citep{JMLR:v13:song12a,kornblith2019similarityneuralnetworkrepresentations,murphy2024correctingbiasedcenteredkernel, cortes2024algorithmslearningkernelsbased} and effective rank \citep{7098875}. Next, we test whether representations inside the tunnel are compatible with the model's pretrained forecasting pathway. Finally, we physically truncate the models and measure the resulting performance-compute frontier to determine how early recoverability translates into inference savings.

Our contributions are:
\begin{enumerate}[leftmargin=*]

\item \textbf{A functional measure of forecasting depth.} We introduce a  sustained forecasting tunnel for TSFMs, which identifies the earliest depth after which all subsequent forecasting readouts remain near the final-block readout. At our primary 5\% tolerance, the tunnel begins before the final block in 40 of 42 model-dataset pairs.

\item \textbf{Recoverability versus geometry.} We show that the onset of the forecasting tunnel is not marked by CKA similarity regimes or effective-rank changes, indicating that geometric redundancy and forecasting redundancy provide different criteria for depth reduction.

\item \textbf{From recoverability to native-pathway compatibility.} We show that intermediate representations can support near-final forecasting readouts while remaining poorly matched to the model's pretrained forecasting pathway. A label-free affine alignment substantially improves this compatibility across all three architectures.

\item \textbf{From recoverability to acceleration.} We physically truncate the models and measure the resulting performance-compute frontier. At matched cut depths, tunnel membership stratifies truncation tolerance in a post hoc comparison, linking the recoverability diagnostic to actual depth reduction. Aligned Chronos-2 remains within a median 3.1\% of the full model while retaining only 4 of 12 blocks and running $2.43\times$ faster; TiRex reaches a $1.71\times$ speedup with an 8.5\% median MASE increase while retaining 7 of 12 blocks, while TimesFM-3 exhibits more modest gains.

\end{enumerate}

The remainder of the paper reviews related work (Section~\ref{sec:related_work}), introduces our methods and hypotheses (Section~\ref{sec:method}), presents the results (Section~\ref{sec:results}), and concludes with limitations (Section~\ref{sec:conclusion}); experimental details and additional results are provided in Appendices~\ref{app:experimental_details} and~\ref{app:additional_results}, respectively.

\section{Related Work}
\label{sec:related_work}

\textbf{Time-series foundation models.}
TSFMs have expanded rapidly in scale and scope, targeting zero-shot, general-purpose forecasting
across heterogeneous domains \citep{ansari2024chronoslearninglanguagetime,
das2024decoderonlyfoundationmodeltimeseries, woo2024unifiedtraininguniversaltime,
ansari2025chronos2univariateuniversalforecasting, auer2025tirexzeroshotforecastinglong}.

\textbf{Layerwise analysis and probing.}
Linear probes measure how readily task-relevant information can be decoded from representations at different depths \citep{alain2018understandingintermediatelayersusing}. Intermediate features can carry information that the final representation does not \citep{evci2022head2toeutilizingintermediaterepresentations}, and in supervised image classifiers linear-probe performance can saturate before the final layer. This is dubbed the ``tunnel effect'' \citep{masarczyk2023tunneleffectbuildingdata}. In language models, the logit lens \citep{nostalgebraist2020logitlens} applies the pretrained output head directly to intermediate states, and the tuned lens \citep{belrose2025elicitinglatentpredictionstransformers} first passes them through a learned affine map. Our forecasting hard cut parallels the logit lens, while our label-free affine adapter has a close analogue in the tuned lens for language models. Representation geometry gives a complementary view: CKA measures similarity between layer representations \citep{kornblith2019similarityneuralnetworkrepresentations, cortes2024algorithmslearningkernelsbased}, and effective rank characterizes their dimensional structure \citep{garrido2023rankmeassessingdownstreamperformance,
skean2025layerlayeruncoveringhidden}. \citet{murphy2024correctingbiasedcenteredkernel} show that standard CKA is inflated when samples are few relative to the representation dimension, even for independent random representations, so we use debiased CKA \citep{JMLR:v13:song12a}.

\textbf{Efficient inference and depth reduction.}
In language models, depth can be removed or bypassed with limited degradation through structured layer dropping \citep{fan2019reducingtransformerdepthdemand}, layer pruning \citep{men2024shortgptlayerslargelanguage, gromov2025unreasonableineffectivenessdeeperlayers}, and early exiting \citep{teerapittayanon2017branchynetfastinferenceearly,
schwartz2020righttooljobmatching, Elhoushi_2024}. These complement broader compression methods such as quantization \citep{frantar2023gptqaccurateposttrainingquantization} and knowledge distillation \citep{hinton2015distillingknowledgeneuralnetwork,
sanh2020distilbertdistilledversionbert}. For time series, \citet{feofanov2026mantisv2closingzeroshotgap} show that simple classifiers trained on intermediate TSFM representations can outperform those trained on the final-layer representation.
\citet{wiliński2025exploringrepresentationsinterventionstime} identify block-like redundancy in TSFMs with CKA and use it to guide pruning.  We instead characterize removable depth directly through forecasting performance.

\section{Method}
\label{sec:method}

We first introduce the notation and four hypotheses that structure our experiments (Section~\ref{sec:intuition}), followed by the layerwise readouts used to measure forecast recoverability (Section~\ref{sec:layerwise_readouts}). We then introduce the hard cut and alignment procedures used to test native-pathway compatibility and acceleration (Section~\ref{sec:adapter}).

\subsection{Setup and hypotheses}
\label{sec:intuition}

A forecasting model maps its input through a sequence of \(L\) blocks. Let \(h^{(0)}\) denote the representation entering the first block, and \(h^{(\ell)}\) the output of block \(\ell\):
\begin{equation}
h^{(\ell)} = F_{\ell}\!\left(h^{(\ell-1)}\right), \qquad \ell = 1,\dots,L,
\label{eq:stage_transform}
\end{equation}
where $F_{\ell}$ is the transformation implemented by block $\ell$ (a transformer block in Chronos-2 and TimesFM-3, and an sLSTM block in TiRex). The model then produces its forecast through its \emph{native forecasting pathway}:
\begin{equation}
\hat{y} = g\!\left(h^{(L)}\right),
\label{eq:native_forecast}
\end{equation}
where $g$ comprises any post-block normalization, the pretrained forecasting head, and the model's inverse target transformation. We index depth only by block outputs, so any final normalization is included in $g$ rather than treated as a separate depth.

\paragraph{Recoverability.}
At each depth $\ell$, we fit a lightweight readout $q_\ell$ on the training split and measure its validation loss:
\begin{equation}
\mathcal{R}_{\ell}^{\mathrm{val}}
= \frac{1}{N_{\mathrm{val}}} \sum_{i=1}^{N_{\mathrm{val}}}
\mathcal{L}\!\left(q_{\ell}\!\left(h_i^{(\ell)}\right), y_i\right),
\label{eq:layerwise_risk}
\end{equation}
where $N_{\mathrm{val}}$ is the number of validation examples, $h_i^{(\ell)}$ the depth-$\ell$ representation of validation example $i$, $y_i$ its target, and $\mathcal{L}$ the per-example forecasting loss. Simple readouts can saturate early in deep image classifiers, a phenomenon known as the tunnel effect \citep{masarczyk2023tunneleffectbuildingdata}; we ask whether the same occurs in time-series forecasting. Rather than taking the first depth that reaches near-final performance, we define a stricter \emph{sustained forecasting tunnel}, beginning at the earliest depth after which every subsequent readout remains within a tolerance $\epsilon$ of the final-block readout:
\begin{equation}
\ell_{\mathrm{tunnel}}(\epsilon)
= \min\Bigl\{ \ell \in \{0,\dots,L\} :
\mathcal{R}_{j}^{\mathrm{val}} \le (1+\epsilon)\,\mathcal{R}_{L}^{\mathrm{val}}
\;\;\forall j \in \{\ell,\dots,L\} \Bigr\}.
\label{eq:tunnel_entrance}
\end{equation}
This sustained criterion prevents an isolated early dip from opening a tunnel. We use $\epsilon=0.05$ as the primary criterion and report $\epsilon\in\{0.02,0.08\}$ as sensitivity analyses. The reference is the \emph{final-block readout}, not the native forecast. Thus, $\ell_{\mathrm{tunnel}}$ marks where later blocks cease to substantially improve forecasting performance recoverable by the readout; it does not indicate where a truncated model would match the full model.

\paragraph{Representation geometry.}
If forecasting performance saturates at $\ell_{\mathrm{tunnel}}$, does representation geometry saturate with it? We track debiased linear CKA between depths \citep{JMLR:v13:song12a,kornblith2019similarityneuralnetworkrepresentations,murphy2024correctingbiasedcenteredkernel,cortes2024algorithmslearningkernelsbased}, which measures similarity between two representations, and entropy effective rank \citep{7098875}, which characterizes how the representation space expands or contracts with depth. If geometry stabilizes at the tunnel entrance, functional and geometric saturation coincide; if it continues to evolve, representation refinement persists after forecasting performance has saturated.

\paragraph{Native-pathway compatibility.}
A layerwise readout establishes only recoverability, while direct truncation requires something stronger: the model's own pathway $g$, trained on $h^{(L)}$, must also operate effectively on $h^{(\ell)}$. Even when $\mathcal{R}_{\ell}^{\mathrm{val}} \approx \mathcal{R}_{L}^{\mathrm{val}}$, directly applying $g$ to $h^{(\ell)}$ can fail because $h^{(\ell)}$ may contain recoverable forecast information while remaining mismatched to the representation expected by $g$. We therefore insert a lightweight alignment map before the pathway:
\begin{equation}
\hat{y}_{\ell} = g\!\left(A_{\ell}\!\left(h^{(\ell)}\right)\right),
\label{eq:truncated_inference}
\end{equation}
where $A_\ell$ is fit to reproduce the full model's forecast (Section~\ref{sec:adapter}). The degradation from directly applying $g$ to $h^{(\ell)}$, relative to the full model, measures the compatibility gap, while the improvement from $A_\ell$ measures how much of that gap can be corrected by a lightweight transformation.

\paragraph{From compatibility to acceleration.}
If $g\!\left(A_{\ell}(h^{(\ell)})\right)$ approaches the full model's forecasting performance, blocks $\ell+1,\dots,L$ can be skipped at inference time. Sweeping $\ell$ then traces a performance-compute frontier. 

We test four hypotheses:

\begin{enumerate}[label=\textbf{(H\arabic*)}, ref=H\arabic*, leftmargin=*]

  \item \label{hyp:recoverability}
  \textbf{Early forecast recoverability.}
  The forecasting tunnel can begin well before the final block, with its entrance varying across models and datasets.

  \item \label{hyp:geometry}
  \textbf{Geometry keeps evolving inside the tunnel.}
  CKA and effective rank continue to change after $\ell_{\mathrm{tunnel}}$ without commensurate gains in forecasting performance accessible to the readout.

  \item \label{hyp:alignment}
  \textbf{Recoverability precedes native-pathway compatibility in TSFMs.}
  At depths where near-final forecasting performance is already recoverable, directly applying $g$ to $h^{(\ell)}$ degrades performance relative to the full model, while lightweight alignment reduces much of this gap.

  \item \label{hyp:acceleration}
  \textbf{Truncate to accelerate.}
  Aligned truncation exposes a performance-compute frontier in which truncation depth trades forecasting performance against inference cost.

\end{enumerate}

\subsection{Layerwise forecasting readouts}
\label{sec:layerwise_readouts}

We probe each model at the token positions read by its native forecasting head. Let $H$ be the
horizon, $d$ the hidden width, $P$ the output patch length, and $Q=9$ the number of quantile levels $\{0.1,\dots,0.9\}$ used across all three models. Chronos-2 forecasts from $K=\lceil H/P\rceil$ forecast-slot tokens, while TiRex forecasts from the last input token of
each of its $K$ inference passes. For both, a shared linear readout
$q_\ell:\mathbb{R}^{d}\to\mathbb{R}^{Q\times P}$ is applied to each forecast-producing state, and the $K$ predicted patches are concatenated. For TimesFM-3, a linear readout
$q_\ell:\mathbb{R}^{d}\to\mathbb{R}^{Q\times H}$ predicts the full horizon from the model's last context token. Readouts are trained with the pinball loss in each model's normalized target space, with weight decay selected per depth on validation (Appendix~\ref{app:readout_training}). Because target spaces differ, absolute losses are compared only within a model.


\subsection{Hard cut and native-pathway alignment}
\label{sec:adapter}

At depth $\ell<L$, the \emph{hard cut} feeds the intermediate representation directly into the pretrained forecasting pathway, $g(h^{(\ell)})$. To test whether a simple representation mismatch explains the resulting degradation, we insert \textit{a label-free affine adapter},
$
A_\ell(h)=h+\Delta_\ell h+b_\ell,
$
initialized at the identity and fitted with the backbone frozen to minimize the squared error between $g(A_\ell(h^{(\ell)}))$ and the full model's own forecast $g(h^{(L)})$. Because TimesFM-3's head is linear, its adapter is fitted in closed form, with a ridge penalty that shrinks $A_\ell$ toward zero (Appendix~\ref{app:alignment}). The adapter needs only unlabeled windows of the target dataset; Appendix~\ref{sec:app_lodo} tests adapters fitted on other datasets instead. We additionally retain the linear readout $q_\ell$ as a supervised reference for how much forecast information is recoverable from the same representation. A nonlinear adapter variant is reported in Appendix~\ref{app:alignment}.


\section{Results}
\label{sec:results}

We first locate the forecasting tunnel (Section~\ref{sec:results_recoverability},
\ref{hyp:recoverability}) and compare it with representation geometry
(Section~\ref{sec:results_geometry}, \ref{hyp:geometry}). We then test whether the pretrained pathway can use early states (Section~\ref{sec:results_alignment}, \ref{hyp:alignment}) and whether aligned truncation accelerates inference (Section~\ref{sec:results_acceleration}, \ref{hyp:acceleration}).

\paragraph{Experimental setup.}
We evaluate Chronos-2, TimesFM-3, and TiRex on 14 datasets spanning five sampling frequencies, from 5-minute to daily: M4 \citep{MAKRIDAKIS202054},
Electricity \citep{electricityloaddiagrams20112014_321, godahewa2021monashtimeseriesforecasting},
Uber TLC \citep{fivethirtyeight_uber_tlc,godahewa2021monashtimeseriesforecasting},
Wind Farms \citep{godahewa2021monashtimeseriesforecasting},
BOOM \citep{cohen2025timedifferentobservabilityperspective},
SG Carpark \citep{qiao2026itstimegenerationtime},
Coastal T-S \citep{qiao2026itstimegenerationtime},
LOOP Seattle \citep{cui2018deep,cui2019traffic},
Electricity (15-minute) \citep{electricityloaddiagrams20112014_321, lai2018modelinglongshorttermtemporal, godahewa2021monashtimeseriesforecasting},
SZ Taxi \citep{jiang2024libcityunifiedlibraryefficient},
London Smart Meters \citep{jeanmidev_smartmeters_london,godahewa2021monashtimeseriesforecasting},
M5 \citep{m5-forecasting-accuracy},
Wiki Daily \citep{ansari2024chronoslearninglanguagetime}, and
Monash Traffic \citep{lai2018modelinglongshorttermtemporal, godahewa2021monashtimeseriesforecasting}. Context length and horizon are fixed at $C=512$ and $H=64$. All three models receive identical windows and splits, and all pretrained backbones remain frozen. Selection uses validation data only, and the test split is evaluated once. We report MASE using each dataset's seasonal period, with 95\% confidence intervals from 5,000 series-level cluster-bootstrap resamples, paired whenever two pathways are compared. We time inference end to end on an A100 over three independent repetitions, using batch size 1 for the main results and batch size 256 for the appendix throughput analysis (Appendix~\ref{app:truncation_results}). Full model configurations, dataset descriptions, split construction, optimization details, and each dataset's documented pretraining status for each model are in Appendix~\ref{app:experimental_details}.

\subsection{Forecasting before the final block}

\label{sec:results_recoverability}
\begin{figure}[t]
  \centering
  \includegraphics[width=0.9\textwidth]{figures/main_representative_q9.pdf}
  \caption{\textbf{Layerwise Q=9 forecasting loss for Electricity, Uber TLC, and London Smart Meters.} Test loss is shown across depth for each architecture. The shaded region begins at the sustained 5\% tunnel entrance, selected using validation loss only; test loss does not affect selection and is shown to assess whether the validation-selected tunnel persists on held-out data. Bands show 95\% series-level cluster-bootstrap intervals ($B=5000$). Complete results are in Appendix~\ref{sec:app_recov_curves}.}
  \label{fig:recov_main}
\end{figure}


We first ask \textit{whether near-final forecasting performance becomes recoverable from intermediate representations}. For convenience, we denote the representation before the first model block as L0 and the outputs of successive blocks as L1, L2, and so on, with the final block at L12 for Chronos-2 and TiRex and L20 for TimesFM-3.
\begin{table}[t]
  \centering
  \caption{\textbf{Sustained 5\% forecasting-tunnel entrances for the Q=9 readout.} L0 denotes the representation before the first model block, while L1, L2, and so on denote the outputs of successive blocks. Entrances are selected from validation loss only. $^\dagger$The final-block readout is within 5\% of the best constant forecast on validation (Appendix~\ref{sec:app_native_reference}).}
  \vspace{4pt}
  \label{tab:tunnel_q9}
  \small
  \setlength{\tabcolsep}{3.5pt}
  \begin{tabular}{@{}lccc@{\qquad}lccc@{}}
  \toprule
  Dataset & Chronos-2 & TimesFM-3 & TiRex & Dataset & Chronos-2 & TimesFM-3 & TiRex \\
  \midrule
  M4 & L8 & L16 & L12 & LOOP Seattle & L9 & L14 & L9 \\
  Electricity & L9 & L15 & L11 & Electricity (15-min) & L10 & L1 & L10 \\
  Uber TLC & L4 & L1 & L10 & SZ Taxi & L2 & L1 & L0$^\dagger$ \\
  Wind Farms & L3 & L0 & L11 & London Smart Meters & L3 & L1 & L8 \\
  SG Carpark & L10 & L12 & L11 & M5 & L0$^\dagger$ & L0$^\dagger$ & L0$^\dagger$ \\
  Coastal T-S & L1 & L5 & L1 & Wiki Daily & L2 & L0 & L1 \\
  BOOM & L3 & L0 & L8 & Monash Traffic & L11 & L12 & L12 \\
  \midrule
  \textbf{Median depth} & \textbf{29\%} & \textbf{5\%} & \textbf{79\%} &
  \textbf{Before final} & \textbf{14/14} & \textbf{14/14} & \textbf{12/14} \\
  \bottomrule
  \end{tabular}
\end{table}


Figure~\ref{fig:recov_main} shows the layerwise $Q=9$ mean quantile loss, while Table~\ref{tab:tunnel_q9} summarizes the sustained 5\% tunnel entrance across all datasets. Among the 42 model-dataset pairs, 40 enter the tunnel before the final block: 14/14 for Chronos-2, 14/14 for TimesFM-3, and 12/14 for TiRex; the two exceptions are TiRex on M4 and Monash Traffic. The depth at which the tunnel begins varies across models and datasets.

The validation-selected tunnel entrances also remain close to final-block performance on held-out test data. Across all 42 model-dataset pairs, the median test-loss increase at the selected entrance is 1.3\% (interquartile range (IQR): $-0.7$\% to 4.3\%), with a maximum increase of 9.8\%. Full held-out results are reported in Appendix~\ref{sec:app_heldout_recoverability}. We additionally repeat the analysis with a substantially lower-capacity, median-only ($Q=1$) readout and obtain highly consistent tunnel entrances relative to the $Q=9$ analysis (Appendix~\ref{sec:app_q1_robustness}). 
We additionally compare our architecture-aware readouts with content-mean pooling and find that the choice of hidden states affects both forecasting loss and the location of the tunnel (Appendix~\ref{sec:app_content_mean}).
Because the tunnel is referenced to the final-block readout, it does not measure closeness to the full model: the final-block readout's validation loss is a median 12.7\% to 20.3\% above the full model's, and a tunnel referenced to the full model exists in only 6 of the 42 pairs (Appendix~\ref{sec:app_native_reference}). Near-native performance is the subject of Sections~\ref{sec:results_alignment} and~\ref{sec:results_acceleration}.


Overall, \textit{near-final forecasting performance can be recovered before the final block}, but how early this occurs depends on the architecture, dataset, and tolerated degradation; with $\epsilon=0.08$, the tunnel begins before the final block in all model-dataset pairs, whereas under the stricter $\epsilon=0.02$ criterion this holds for 36/42 pairs (Appendix~\ref{sec:app_tunnel_sensitivity}). We use the 5\% criterion for our main analysis and next ask \textit{whether representation geometry also saturates once near-final forecasting performance becomes recoverable}. 

\subsection{Geometry and forecast recoverability}
\label{sec:results_geometry}

\begin{figure}[t]
  \centering
  \includegraphics[width=0.9\textwidth]{figures/geometry_main.pdf}
  \caption{\textbf{Representation geometry and forecast recoverability.}
    Effective rank and debiased linear CKA between each depth and the final block, computed on training representations, for SG Carpark (top) and London Smart Meters (bottom). Columns correspond to Chronos-2, TimesFM-3, and TiRex. Shading marks the validation-selected 5\% forecasting tunnel.}
  \label{fig:geometry_main}
\end{figure}

Figure~\ref{fig:geometry_main} shows that the forecasting tunnel does not consistently align with changes in representation geometry. On SG Carpark, the Chronos-2 tunnel begins several blocks after both the effective-rank peak and a major CKA transition; for TimesFM-3, effective rank continues to increase after the tunnel entrance; and for TiRex, substantial CKA evolution continues after the tunnel has already begun. Similarly, for London Smart Meters, the tunnel begins at L3, L1, and L8 for Chronos-2, TimesFM-3, and TiRex, respectively, while effective rank and similarity to the final representation continue to evolve at later depths. 
These patterns persist under the stricter 2\% tolerance (Appendix~\ref{app:geometry}).

Taken together, these results show that \textit{forecasting performance can saturate while geometry keeps evolving}. Moreover, geometry alone  neither locates the tunnel nor defines a unique truncation depth: the effective-rank peak identifies a single layer with no direct link to forecast quality, while CKA can reveal several blocks of mutually similar layers, each providing a possible cut point (see Appendix~\ref{app:geometry} for Chronos-2). It does, however, flag a potential challenge for truncation: at the tunnel entrance, CKA with the final block is below 0.5 in 55\% of model-dataset pairs, so \textit{the pretrained head may receive a substantially different representation}.

\subsection{Native-pathway compatibility across depth}
\label{sec:results_alignment}

Using the hard cut and label-free affine adapter introduced in Section~\ref{sec:adapter}, Figure~\ref{fig:h3_main} shows a clear separation between forecast recoverability and native-pathway compatibility. On Electricity, the hard cut never stays within 5\% of the full model before the final block, while the affine adapter closes a median 73\%, over 99\%, and 82\% of the hard-cut gap across intermediate depths for Chronos-2, TimesFM-3, and TiRex, respectively; the particularly large value for TimesFM-3 reflects its extremely poor hard-cut baseline. For all three models, the aligned pathway outperforms both the hard cut and the supervised readout $q_\ell$ used in Section~\ref{sec:results_recoverability} at every depth reported in Table~\ref{tab:h3_depth}. For TimesFM-3, this requires regularizing its closed-form adapter toward zero rather than toward the hard cut: anchored at the hard cut, which is 35 to 160 times worse than the full model, the adapter overfits its training windows and degrades by $+48\%$ to $+70\%$ at 25\% to 75\% of depth (Appendix~\ref{sec:app_timesfm3_ridge}). A nonlinear bottleneck adapter changes little relative to the affine adapter for Chronos-2 and TiRex, with a median difference of only $-0.2$ percentage points in test MASE, whereas it performs substantially worse for TimesFM-3 (Appendix~\ref{app:compatibility}).

We conclude that \textit{recoverability precedes native-pathway compatibility in TSFMs}, and that lightweight alignment substantially improves compatibility across all three architectures. The extent of recovery is architecture-dependent: for Chronos-2 and TiRex, a depth-specific affine map can bring intermediate states close to native performance several blocks before the final one, while for TimesFM-3 the aligned pathway stays about 10\% above the full model from L1 to L14 and comes within 5\% of it on most datasets only at L18.
 
\begin{figure}[t]
  \centering
  \includegraphics[width=0.9\textwidth]{figures/h3_main_electricity_rev.pdf}
  \caption{\textbf{Native-pathway compatibility across depth on Electricity.} Test MASE increase over the full model for the hard cut, affine-aligned pathway, and supervised probe reference. Shading marks the forecasting tunnel; bands show 95\% cluster-bootstrap intervals. The y-axis is linear to $+20\%$ and log-scaled above; triangles mark clipped values. The TimesFM-3 adapter is regularized toward zero (Appendix~\ref{sec:app_timesfm3_ridge}). Complete results are in Appendix~\ref{app:compatibility}.}
  \label{fig:h3_main}
\end{figure}



\begin{table}[t]
\centering
\caption{\textbf{Native-pathway compatibility across depth.} Median test MASE increase over the full model across 14 datasets at 25\%, 50\%, and 75\% of network depth, corresponding to L3/L6/L9 for Chronos-2 and TiRex and L5/L10/L15 for TimesFM-3. $\times$ denotes ratios to the full-model MASE. TimesFM-3 uses the zero-anchored adapter; anchored at the hard cut it gives $+69.8\%$, $+67.3\%$, and $+47.9\%$ (Appendix~\ref{sec:app_timesfm3_ridge}).}
\vspace{4pt}
\label{tab:h3_depth}
\small
\setlength{\tabcolsep}{4pt}
\begin{tabular}{@{}lccc ccc ccc@{}}
\toprule
& \multicolumn{3}{c}{Hard cut}
& \multicolumn{3}{c}{Aligned pathway}
& \multicolumn{3}{c}{Probe reference} \\
\cmidrule(lr){2-4}
\cmidrule(lr){5-7}
\cmidrule(l){8-10}
Model
& 25\% & 50\% & 75\%
& 25\% & 50\% & 75\%
& 25\% & 50\% & 75\% \\
\midrule
Chronos-2
& $+44.2\%$ & $+20.1\%$ & $+17.1\%$
& $+4.4\%$ & $+2.3\%$ & $+2.2\%$
& $+13.9\%$ & $+11.8\%$ & $+11.6\%$ \\

TimesFM-3
& $53\times$ & $121\times$ & $160\times$
& $+10.6\%$ & $+11.3\%$ & $+8.5\%$
& $+16.8\%$ & $+17.0\%$ & $+14.8\%$ \\

TiRex
& $+121.3\%$ & $+104.7\%$ & $+80.0\%$
& $+13.4\%$ & $+10.6\%$ & $+6.2\%$
& $+14.5\%$ & $+13.2\%$ & $+12.8\%$ \\
\bottomrule
\end{tabular}
\end{table}

\subsection{Truncate to accelerate}
\label{sec:results_acceleration}

The forecasting tunnel establishes that near-final forecast information often becomes recoverable well before the final block, while Section~\ref{sec:results_alignment} shows that alignment can make such intermediate states substantially more compatible with the native pathway. We now ask whether this early recoverability can be converted into actual inference savings.

\textit{How much does it cost to remove blocks?} Figure~\ref{fig:h4_frontier} shows the resulting \textit{performance-compute frontier}, obtained by physically truncating each model, verifying agreement with the offline aligned pathway in mean test MASE to within $10^{-4}$ relative error (Appendix~\ref{app:latency}), and measuring end-to-end speedup over the full model. Aligned Chronos-2 remains within a median 3.1\% of the full model from L11 through L4, where it runs $2.43\times$ faster. TiRex exhibits a more gradual tradeoff, reaching $1.71\times$ speedup at an 8.5\% median MASE increase when truncated after L7. TimesFM-3 provides more modest gains at batch size 1: even removing all blocks reaches only $1.38\times$, indicating that a substantial fraction of its end-to-end runtime lies outside the Transformer blocks (Table~\ref{tab:h4_operating_points}); its aligned pathway stays about 10\% above the full model from L1 to L14, with accuracy computed offline (Appendix~\ref{sec:app_timesfm3_ridge}). Chronos-2-small \citep{ansari2025chronos2univariateuniversalforecasting} provides a strong separately pretrained baseline, running $1.80\times$ faster at only $+0.8\%$ median MASE; at a comparable speed, aligned truncation after L6 reaches $1.78\times$ with a $+2.3\%$ increase. At batch size 256 its advantage is larger ($3.28\times$, against $3.10\times$ at $+4.4\%$ for aligned truncation after L3). Truncation is therefore most useful beyond Chronos-2-small's speed at batch size 1, for models without a smaller pretrained version, such as TiRex and TimesFM-3, and for obtaining several speed-performance tradeoffs from a single pretrained checkpoint without training a separate model. Complete layer-by-layer performance and speedup results are reported in Appendix~\ref{app:truncation_results}.

\textit{Does one adapter serve unseen datasets?} The adapters above are fitted per dataset on unlabeled windows; for Chronos-2, this takes one forward pass of the full model over 1,394 windows and about one minute of fitting on one A100. An adapter fitted only on the other datasets, without any data from the target, raises the median MASE increase from $+3.1\%$ to $+12.7\%$ after L4 and from $+2.2\%$ to $+4.0\%$ after L9, compared with $+37.0\%$ and $+17.1\%$ for the hard cut, and improves on the hard cut on 13 of 14 datasets at every depth; the exception is LOOP Seattle, our only 5-minute dataset (Appendix~\ref{sec:app_lodo}). Pooling more than about four source datasets brings little further gain (Appendix~\ref{sec:app_source_diversity}). Truncation without target-specific calibration is therefore viable at mild truncation, whereas aggressive truncation still benefits from calibration on the target dataset.

\textit{Does the forecasting tunnel identify where truncation is well tolerated?}
Because inference speed depends only on the cut depth, we compare datasets at the same cut according to whether that depth lies inside or before their validation-selected tunnel. Despite being defined only from layerwise readouts, without evaluating a truncated model, the tunnel stratifies truncation performance in this post hoc comparison. After L4, aligned Chronos-2 incurs a median $+2.2\%$ MASE increase on datasets for which L4 lies inside the tunnel, compared with $+14.2\%$ when it lies before the tunnel; after L6, the corresponding TiRex increases are $+4.4\%$ and $+12.3\%$. Truncation is cheaper inside the tunnel at every depth where both groups contain at least two datasets, for all three models (Table~\ref{tab:h4_tunnel_strata}), although some groups contain only two datasets and datasets whose tunnel begins early may be easier to truncate for other reasons. For TimesFM-3, the aligned pathway remains about 10\% above the full model inside the tunnel from L1 to L14, even though its tunnel begins at a median 5\% of depth, so early recoverability alone does not ensure native-pathway compatibility. Cutting each dataset at its own tunnel entrance, a rule we fixed before computing the truncation results, keeps the aligned model within 5\% of the full model on 7 of 14 datasets for Chronos-2 and 5 of 12 for TiRex, and roughly halves the worst case of a single fixed cut at a similar speed (Appendix~\ref{sec:app_entrance_rule}). Thus, the forecasting tunnel identifies datasets that better tolerate a given cut, even though its entrance need not define the best truncation depth.

\begin{table}[t]
\centering
\caption{\textbf{Operating points of aligned truncation.}
Median test MASE increase across 14 datasets and end-to-end speedup at batch size 1. For TimesFM-3 at L19, the adapter overhead slightly exceeds the saving from removing one block. TimesFM-3 uses the zero-anchored adapter, whose accuracy is computed offline (Appendix~\ref{sec:app_timesfm3_ridge}). Full layer-by-layer results are reported in Appendix~\ref{app:truncation_results}.}
\vspace{4pt}
\label{tab:h4_operating_points}
\small
\setlength{\tabcolsep}{4pt}
\begin{tabular}{@{}lrrr rrr rrr@{}}
\toprule
Model
& Layer & $\Delta$MASE & Speedup
& Layer & $\Delta$MASE & Speedup
& Layer & $\Delta$MASE & Speedup \\
\midrule

Chronos-2
& L9 & $+2.2\%$ & $1.28\times$
& L6 & $+2.3\%$ & $1.78\times$
& L4 & $+3.1\%$ & $2.43\times$ \\

TimesFM-3
& L19 & $+1.3\%$ & $0.99\times$
& L10 & $+11.3\%$ & $1.19\times$
& L0 & $+24.3\%$ & $1.38\times$ \\

TiRex
& L11 & $+1.8\%$ & $1.09\times$
& L7 & $+8.5\%$ & $1.71\times$
& L3 & $+13.4\%$ & $3.93\times$ \\

\bottomrule
\end{tabular}
\end{table}

\begin{figure}[t]
  \centering
  \includegraphics[width=0.9\textwidth]{figures/h4_frontier_rev.pdf}
  \caption{\textbf{Performance-compute frontier of truncation.} Median test MASE increase across 14 datasets versus end-to-end speedup at batch size 1. Points denote aligned cut depths; the star is the full model and the diamond is Chronos-2-small; bands show interquartile ranges. TimesFM-3 uses the zero-anchored adapter (Appendix~\ref{sec:app_timesfm3_ridge}).}
  \label{fig:h4_frontier}
\end{figure}


\section{Conclusion and Limitations}
\label{sec:conclusion}

Together, these results \textit{complete the progression from layerwise analysis to acceleration}. Rather than relying on geometric similarity alone to identify removable depth in time-series foundation models, we introduce the forecasting tunnel, a direct task-level criterion based on forecasting performance and a single tolerance parameter. It shows that near-final forecast information can become recoverable well before the final block even while representation geometry continues to evolve. We then ask whether these early representations can be reused by the model's pretrained forecasting pathway: direct reuse often fails, but lightweight alignment substantially improves compatibility. Finally, physical truncation converts this compatibility into a measured performance-compute frontier, \textit{allowing different speed-performance operating points from a single pretrained checkpoint}.

Across three state-of-the-art TSFMs and 14 datasets spanning five sampling frequencies, the forecasting tunnel begins before the final block in 40 of 42 model-dataset pairs at our primary 5\% tolerance, and remains pre-final in 36 of 42 pairs even under the stricter 2\% tolerance. Representation geometry nevertheless continues to evolve inside the tunnel, showing that forecasting saturation need not coincide with geometric stabilization. Direct reuse of intermediate states through the pretrained forecasting pathway rarely succeeds, whereas lightweight affine alignment substantially improves compatibility across all three architectures. Physical truncation then turns this improved compatibility into measurable inference savings: Chronos-2 remains within a median 3.1\% of the full model after retaining only 4 of its 12 blocks while running $2.43\times$ faster, and TiRex reaches $1.71\times$ speedup with an 8.5\% median MASE increase while retaining 7 of 12 blocks. Importantly, at matched cut depths, datasets for which the cut lies inside the forecasting tunnel generally tolerate truncation better, directly linking early recoverability to practical depth reduction. TimesFM-3 provides the complementary case: its tunnel begins at a median 5\% of depth, but its aligned pathway stays about 10\% above the full model up to L14 and comes within 5\% of it on most datasets only at L18, showing that recoverability alone does not guarantee native-pathway compatibility. Its batch-size-1 speedups are also modest, with truncation before the first Transformer block reaching only $1.38\times$, consistent with a substantial fraction of end-to-end runtime lying outside the Transformer blocks. Chronos-2-small remains a strong separately pretrained baseline, while truncation provides multiple post-hoc operating points from a single pretrained checkpoint.

Our readouts and adapters are deliberately simple, making our estimates of recoverability and compatibility conservative with respect to more expressive mappings. For TimesFM-3, the affine adapter still leaves a gap of about 10\% to the full model at intermediate depths, fails on BOOM and Wind Farms, where none of our affine fits reproduces the full model's forecasts on validation windows, and is evaluated offline (Appendix~\ref{sec:app_timesfm3_ridge}); richer alignment, for example over the full token sequence, may close more of the gap. Adapters are fitted per dataset on unlabeled windows. An adapter fitted on other datasets transfers well enough for mild truncation, but not for aggressive truncation, and fails on LOOP Seattle (Appendices~\ref{sec:app_lodo} and~\ref{sec:app_source_diversity}), so a single truncated model at each depth for all datasets remains limited to mild truncation. We also truncate contiguous prefixes only; for Chronos-2 at mild compression, skipping interior blocks while keeping the final ones is better (Appendix~\ref{sec:app_noncontig}). With 14 datasets, we can establish an association between forecast recoverability and truncation tolerance but not determine how reliably it generalizes or supports cut-depth selection. Larger benchmarks such as GIFT-Eval \citep{aksu2024giftevalbenchmarkgeneraltime} and fev-bench \citep{shchur2026fevbenchrealisticbenchmarktime} could test whether the forecasting tunnel can serve as a reliable indicator of truncation tolerance and guide cut-depth selection. Several datasets in our study are documented pretraining data for at least one model (Appendix~\ref{app:provenance}), so frozen-backbone experiments cannot fully separate generalizable structure from effects of pretraining exposure; extending the analysis to fev-bench, which is held out from the pretraining of Chronos-2 and TimesFM-3, would provide a cleaner test. Finally, we study univariate forecasting at a single context length and horizon, with speed measured on one GPU type; multivariate and covariate-aware forecasting, other horizons and context lengths, and additional hardware remain important directions for future work.

Together, these findings mark a step toward TSFMs that are both easier to understand internally and more efficient to deploy, by connecting where forecast information becomes recoverable to where computation can be removed.

\subsection*{AI use statement}

In this work, we used generative AI tools for writing assistance, such as improving word choice, correcting grammar, and improving the flow and transitions between sentences. Generative AI tools also assisted in drafting the Appendix sections describing the latency protocol and compute accounting; the underlying experimental procedures, measurements, and reported values were determined and verified by the authors, and all generated text was reviewed and revised by the authors. We also used generative AI tools for coding, with their behavior constrained by author-written instruction files. All AI-generated code was reviewed and debugged by the authors and validated through tests written both by the authors and with AI assistance. Additionally, generative AI tools were used to implement plotting and table-generation code for the paper, under the authors' review. Generative AI tools were also used to assist in writing scripts for collecting latency and compute information, such as device metadata, runtime records, and speedup measurements; all such code was reviewed by the authors. During the review period, generative AI tools assisted in analyzing the fits of the original TimesFM-3 adapter, in implementing and running the additional experiments, and in drafting the appendices added in response to the reviews (Appendices~\ref{sec:app_native_reference}, \ref{sec:app_timesfm3_ridge}, and~\ref{sec:app_entrance_rule} to~\ref{sec:app_heldout_pairs}) and the corresponding revisions of the main text; all such text was reviewed and revised by the authors.

We have not used generative AI tools to generate research ideas, conduct the literature review, design methods, select time-series foundation models or datasets, design experiments, or draft sections of the main paper as originally submitted. We take responsibility for the final content of this work, including text, claims, or artifacts produced with the aid of generative AI.

\subsection*{Ethics statement}

This work uses three publicly released pretrained forecasting models and 14 public datasets and involves no human subjects. Most datasets contain aggregated counts, sensor measurements, or system metrics, including traffic, parking, ride demand, wind power, oceanographic, and cloud-observability data. The two electricity-consumption datasets contain anonymized readings from individual clients or households; we use them only in the form released by their providers and make no attempt at re-identification. Our methods analyze and truncate existing models rather than introduce new capabilities, and can reduce the computational cost of inference. Because truncation trades forecasting performance against speed, with the attainable tradeoff varying across models and datasets, truncated models should be validated on the target data before use in consequential decisions.


\subsection*{Reproducibility statement}

All three models are public checkpoints, used frozen; the exact checkpoints and package versions are listed in Appendix~\ref{app:models}. All 14 datasets are publicly available, and Appendix~\ref{app:datasets} gives their sources and every preprocessing step. The construction of windows and splits, including the fixed random seed used for random choices, is described in Appendix~\ref{app:splits}; all three models are evaluated on identical windows, checked against a stored reference for each dataset. The target spaces, readout training, alignment procedure, and metrics are specified in Appendices~\ref{app:targets}, \ref{app:readout_training}, \ref{app:alignment}, and~\ref{app:metrics}, including all hyperparameter grids and selection rules.

The tunnel definition and all selection rules use validation data only and were fixed before the corresponding test results were computed. The comparison of truncation inside and before the tunnel (Section~\ref{sec:results_acceleration}) is a post hoc analysis. The zero-anchored TimesFM-3 adapter was introduced after the test results of the original, identity-anchored adapter had been computed; both are reported (Appendix~\ref{sec:app_timesfm3_ridge}). The analyses added during the review period (Appendices~\ref{sec:app_content_mean}, \ref{sec:app_native_reference}, and~\ref{sec:app_lodo} to~\ref{sec:app_heldout_pairs}) are post hoc, apart from the test of cutting at the tunnel entrance (Appendix~\ref{sec:app_entrance_rule}), which was specified before the truncation results; their own targets, block selections, and selection rules were fixed before their results were computed. A pre-specified per-window equivalence check between the offline aligned pathway and the physically truncated implementation was not satisfied at every operating point; however, mean test MASE agrees to within $10^{-4}$ relative error. Appendix~\ref{app:latency} reports this check together with the latency protocol. Appendix~\ref{app:compute} lists the hardware, software, and compute.




\bibliography{iclr2027_conference}
\bibliographystyle{iclr2027_conference}

\clearpage
\appendix
\section{Experimental Details}
\label{app:experimental_details}


\subsection{Models}
\label{app:models}

All three models are used as released, with frozen float32 parameters, and run in univariate mode. Table~\ref{tab:models} summarizes the architecture and readout positions of each model. For every model, depth is indexed as $\mathrm{L}0,\dots,\mathrm{L}L$: $\mathrm{L}0$ denotes the representation entering the first block, and $\mathrm{L}\ell$ the output of block $\ell$. Any normalization applied after the final block belongs to the native forecasting pathway $g$ and does not define an additional depth. We still probe its output as a diagnostic, but exclude it from the forecasting-tunnel definition and the main CKA analysis. All models are evaluated at the same nine quantile levels $\{0.1,\dots,0.9\}$, verified against each model's native implementation.

\begin{table}[h]
  \centering
  \small
  \caption{\textbf{Models and readout positions.} $L$: number of blocks. $d$: hidden width.
  $K$: forecast-producing states per window at $H=64$, each covering $P$ horizon steps.}
  \label{tab:models}
  \begin{tabularx}{\textwidth}{@{}lXXX@{}}
    \toprule
    & Chronos-2 & TimesFM-3 & TiRex \\
    \midrule
    Checkpoint & \texttt{amazon/chronos-2} & \texttt{google/timesfm-3.0} & \texttt{NX-AI/TiRex} \\
     &  & \texttt{pytorch} &  \\
    Package & \texttt{chronos-forecasting} 2.3.1 & \texttt{timesfm} 3.0.2 & \texttt{tirex-ts} 1.4.2 \\
    Backbone & encoder-only Transformer & decoder-only Transformer & xLSTM (sLSTM blocks) \\
    $L$ / $d$ & 12 / 768 & 20 / 1280 & 12 / 512 \\
    Input patch & 16 & 32 & 32 \\
    Readout states & $K{=}4$ forecast-slot tokens & last context token (index 15) & last input token (index 63) of each of $K{=}2$ passes \\
    $P$ per state & 16 & 64 & 32 \\
    $\mathrm{L}0$ & tokens entering block 1 & output of the input residual block & output of the input patch embedding \\
    Head input & T5-style RMS normalization of L12 & L20 & final RMSNorm of L12 \\
    Native head & residual MLP block & linear, $1280\to 64\cdot 9$ & residual MLP block, $512\to 2048\to 32\cdot 9$ \\
    Target space & standardize, then arcsinh & detrend, then RevIN & per-pass location/scale \\
    Readout $q_\ell$ & $\mathbb{R}^{768}\!\to\!\mathbb{R}^{9\times 16}$, shared & $\mathbb{R}^{1280}\!\to\!\mathbb{R}^{9\times 64}$ & $\mathbb{R}^{512}\!\to\!\mathbb{R}^{9\times 32}$, shared \\
    \bottomrule
  \end{tabularx}
\end{table}

\paragraph{Chronos-2.} We run one encoder pass with $K=\lceil H/16\rceil=4$ forecast slots, the model's native layout for $H=64$, and read the four slot states. Because encoder attention is bidirectional, the slot states depend on $K$, so all Chronos-2 features come from this $K=4$ pass. The native head predicts 21 quantiles, and we use the nine that match the other two models. Its final normalization is a T5-style layer norm, i.e.\ an RMS normalization without mean subtraction or bias.

\paragraph{TimesFM-3.} With $C=512$ and input patch 32, the context occupies 16 tokens. The native head reads only the last one (index 15), and one output patch of 64 steps covers the whole horizon, so $K=1$. Its forecast is decoded by inverting RevIN with that token's statistics, clipping, and adding back the context trend.

\paragraph{TiRex.} TiRex left-pads every context with missing values to its training length of 2048, so $C=512$ becomes 48 masked and 16 real patches. We use the package's default inference path, which produces $H=64$ in two passes: the second pass appends one patch of missing values, rather than predicted values, to the context. Each pass forecasts 32 steps from token 63, with its own location and scale. We use the package's PyTorch backend throughout. Its sLSTM recurrence runs in bfloat16 with pointwise operations in float32, causing hidden states to vary slightly with GPU batch size; we therefore extract all features at a fixed batch size of 256. The second pass's L0 state is identical across windows because the appended patch is fully masked, so it carries no window-specific information.

\paragraph{Native-pathway check.} For every model and dataset, applying the native pathway to the hooked final-block state reproduces the package's own forecast up to floating-point accumulation error.

\subsection{Datasets}
\label{app:datasets}

Table~\ref{tab:datasets} lists the 14 datasets. Every dataset is used univariately: a
multivariate dataset contributes each selected \emph{variate}, i.e.\ each channel, as a separate
series. Several datasets come from three public collections: the Monash Time Series
Forecasting Archive \citep{godahewa2021monashtimeseriesforecasting}; the UCI Machine Learning
Repository, a public dataset archive hosted by the University of California, Irvine; and TIME, a
recent forecasting benchmark \citep{qiao2026itstimegenerationtime}. Series lengths are reported
before windowing. London Smart Meters, M5, and Wiki Daily are additionally reduced to a
fixed-seed random sample of 1,394 eligible series (marked $^\dagger$ in
Table~\ref{tab:datasets}; Appendix~\ref{app:splits}).

\begin{table}[h]
  \centering
  \small
  \caption{\textbf{Datasets.} $m$: seasonal period used by MASE. Series: number of series after
  dataset-specific construction or filtering but before the fixed-seed subsampling marked by
  $^\dagger$. Lengths are in time steps at the frequency used.}
  \label{tab:datasets}
  \vspace{4pt}
  \begin{tabular}{@{}llrrl@{}}
    \toprule
    Dataset & Freq. & $m$ & Series & Length (min / max) \\
    \midrule
    M4 & 1H & 24 & 414 & 748 / 1,008 \\
    Electricity & 1H & 24 & 321 & 26,304 \\
    Uber TLC & 1H & 24 & 262 & 4,344 \\
    Wind Farms & 1H & 24 & 337 & 1,715 / 8,784 \\
    BOOM & 1H & 24 & 356 & 664 / 13,159 \\
    SG Carpark & 1H (from 15min) & 24 & 354 & $\approx$3,623 \\
    Coastal T-S & 1H & 24 & 48 (24 stations) & 2,733 / 8,784 \\
    LOOP Seattle & 5min & 288 & 323 & 105,120 \\
    Electricity (15-minute) & 15min & 96 & 370 & 16,032 / 140,256 \\
    SZ Taxi & 15min & 96 & 156 & 2,976 \\
    London Smart Meters & 30min & 48 & 5,560$^\dagger$ & 288 / 39,648 \\
    M5 & 1D & 7 & 30,490$^\dagger$ & 124 / 1,969 \\
    Wiki Daily & 1D & 7 & 100,000$^\dagger$ & 2,741 \\
    Monash Traffic & 1H & 24 & 862 & 17,544 \\
    \bottomrule
  \end{tabular}
\end{table}

\paragraph{M4.}
The 414 hourly series of the M4 forecasting competition, anonymized and drawn from several
domains \citep{MAKRIDAKIS202054}.

\paragraph{Electricity.}
Hourly electricity load from 321 clients (2012 to 2014). It is derived from the UCI dataset \citep{electricityloaddiagrams20112014_321}, which records 370 clients every 15 minutes. The Monash archive distributes an hourly aggregate of 321 of these clients, the version used by
\citet{lai2018modelinglongshorttermtemporal}
\citep{godahewa2021monashtimeseriesforecasting}.

\paragraph{Uber TLC.}
Hourly Uber pickup counts across New York City locations, derived from data released by the
New York City Taxi and Limousine Commission
\citep{fivethirtyeight_uber_tlc,godahewa2021monashtimeseriesforecasting}.

\paragraph{Wind Farms.}
Hourly power production from 337 Australian wind farms
\citep{godahewa2021monashtimeseriesforecasting}. 

\paragraph{BOOM.}
Monitoring metrics from Datadog production systems
\citep{cohen2025timedifferentobservabilityperspective}. We use the 378 queries
recorded hourly and, for each query, select the first variate with at most 20\% missing values and at least one complete window, retaining 356 queries.

\paragraph{SG Carpark.}
Available-lot counts from 354 public carparks in Singapore (2025), from TIME
\citep{qiao2026itstimegenerationtime}. We aggregate the 15-minute data to hourly means,
requiring at least three of the four samples in an hour and never filling missing values.

\paragraph{Coastal T-S.}
Hourly measurements from 24 moored instruments, i.e.\ sensors anchored at fixed points, along
the Australian coast, from TIME \citep{qiao2026itstimegenerationtime}. We use temperature and salinity, giving 48 series. The temperature and salinity series from
each station form one bootstrap cluster, leaving 24 independent clusters and producing the widest confidence intervals in the benchmark.

\paragraph{LOOP Seattle.}
Five-minute freeway speeds from 323 sensor stations in the Seattle area during 2015 
\citep{cui2018deep,cui2019traffic}.

\paragraph{Electricity (15-minute).}
The UCI dataset itself \citep{electricityloaddiagrams20112014_321}: all 370 clients at the native
15-minute resolution (2011 to 2014)
\citep{godahewa2021monashtimeseriesforecasting}. It shares its source with Electricity and serves as a \emph{frequency control}: it keeps the domain and most meters fixed while changing the sampling rate. With $C=512$, its context spans 5.3 days, compared with 21 days for hourly Electricity.

\paragraph{SZ Taxi.}
Fifteen-minute road speeds on 156 roads in Shenzhen during January 2015, computed from taxi GPS trajectories. We use the version distributed by LibCity \citep{jiang2024libcityunifiedlibraryefficient}.

\paragraph{London Smart Meters.}
Half-hourly electricity consumption from 5,560 London households (November 2011 to February
2014), collected by the electricity distributor UK Power Networks in its Low Carbon London
smart-meter trial
\citep{jeanmidev_smartmeters_london,godahewa2021monashtimeseriesforecasting}. We use the Monash
version that retains missing values.

\paragraph{M5.}
Daily unit sales from 30,490 Walmart item-store series (2011 to 2016) from the M5 forecasting
competition \citep{m5-forecasting-accuracy}, including the competition's final evaluation
period.

\paragraph{Wiki Daily.}
Daily page views from 100,000 English Wikipedia articles (2015 to 2022), compiled by
\citet{ansari2024chronoslearninglanguagetime}.

\paragraph{Monash Traffic.}
Hourly road occupancy rates from 862 sensors on San Francisco Bay Area freeways (2015 to 2016). We use the version distributed by the Monash archive \citep{godahewa2021monashtimeseriesforecasting}, which was also used by \citet{lai2018modelinglongshorttermtemporal}.

\subsection{Pretraining provenance}
\label{app:provenance}

Because all three models were pretrained on large public corpora, some evaluation datasets may
overlap their pretraining data. For every model-dataset pair, we record what the model's own
documentation establishes using three categories:
\begin{itemize}[leftmargin=*]
  \item \textbf{Listed}: the dataset appears in the documented pretraining corpus.
  \item \textbf{Held out}: the documentation explicitly excludes the dataset, or a benchmark
  containing it, from the relevant real-data pretraining corpus.
  \item \textbf{Not listed}: the dataset appears in neither category. This means only that the
  available documentation does not identify it as training data; it does \emph{not} establish
  that the model never saw it.
\end{itemize}
For each cell, we additionally distinguish evidence for the exact dataset from evidence covering
only the same source at a different resolution or in a different form, marked \emph{source}.

\paragraph{Sources.}
\emph{Chronos-2} lists its real univariate pretraining datasets in Table~6 of
\citet{ansari2025chronos2univariateuniversalforecasting}. Its authors also state that no dataset or task from fev-bench \citep{shchur2026fevbenchrealisticbenchmarktime} or Chronos Benchmark~II
\citep{ansari2024chronoslearninglanguagetime} was used for training.
\emph{TimesFM-3} has no dedicated paper; its model card describes its training corpus as the
GIFT-Eval pretraining collection \citep{aksu2024giftevalbenchmarkgeneraltime}, excluding
datasets that overlap fev-bench, together with Wikipedia page views through November 2023,
Google Trends queries through the end of 2022, and unspecified synthetic and augmented data.
\emph{TiRex} documents its real pretraining data in Tables~5 and~6 of
\citet{auer2025tirexzeroshotforecastinglong}: Chronos training datasets and a subset of the
GIFT-Eval pretraining collection, supplemented with synthetic data. Its authors state that the
training data do not overlap Chronos Benchmark~II. Benchmark membership was checked against
the published task lists: Chronos Benchmark~II contains M5 and Traffic, while fev-bench
contains LOOP Seattle (5-minute), SZ Taxi (15-minute), M5, and 15 tasks derived from BOOM.

\begin{table}[h]
  \centering
  \small
  \caption{\textbf{Pretraining provenance.} Listed: in the documented pretraining corpus.
  Held out: explicitly excluded through a dataset or benchmark-level holdout. Not listed:
  identified by neither, which does not imply unseen. \emph{source}: evidence covers the same
  source at a different resolution or in a different form. Evidence: [C2] Chronos-2 Table~6;
  [BII] Chronos Benchmark~II exclusion; [fev] fev-bench exclusion; [GE] GIFT-Eval pretraining
  collection, via the TimesFM-3 model card; [WP] Wikipedia page views, via the TimesFM-3 model
  card; [T5] TiRex Table~5. $^{\ast}$ fev-bench contains tasks derived from BOOM; see Notes.}
  \label{tab:provenance}
  \begin{tabular}{@{}llll@{}}
    \toprule
    Dataset & Chronos-2 & TimesFM-3 & TiRex \\
    \midrule
    M4 & listed [C2] & not listed & listed [T5] \\
    Electricity & listed [C2] & not listed & listed [T5] \\
    Uber TLC & listed, \emph{source} [C2] & listed [GE] & listed [T5] \\
    Wind Farms & listed [C2] & listed, \emph{source} [GE] & listed [T5] \\
    BOOM & not listed$^{\ast}$ & not listed$^{\ast}$ & not listed \\
    SG Carpark & not listed & not listed & not listed \\
    Coastal T-S & not listed & not listed & not listed \\
    LOOP Seattle & held out [fev] & held out [fev] & not listed \\
    Electricity (15-minute) & listed [C2] & not listed & listed [T5] \\
    SZ Taxi & held out [fev] & held out [fev] & not listed \\
    London Smart Meters & listed [C2] & listed [GE] & listed [T5] \\
    M5 & held out [BII, fev] & held out [fev] & held out [BII] \\
    Wiki Daily & listed [C2] & listed, \emph{source} [WP] & listed [T5] \\
    Monash Traffic & held out [BII] & listed [GE] & held out [BII] \\
    \midrule
    Listed / held out / not listed & 7 / 4 / 3 & 5 / 3 / 6 & 7 / 2 / 5 \\
    \bottomrule
  \end{tabular}
\end{table}

\paragraph{Notes.}
\begin{itemize}[leftmargin=*]
  \item \emph{BOOM}$^{\ast}$: fev-bench contains 15 tasks derived from BOOM. Chronos-2
  excludes fev-bench from training, and TimesFM-3 excludes fev-bench overlaps from its
  GIFT-Eval pretraining component. We could not establish whether those tasks share series
  with the 356 hourly queries used here, so we conservatively record BOOM as not listed.
  \item \emph{Wind Farms, TimesFM-3}: the GIFT-Eval pretraining collection contains the same
  wind farms at minute resolution; we evaluate hourly values.
  \item \emph{Wiki Daily, TimesFM-3}: the model card names Wikipedia page views through
  November 2023 as a training source, which covers the time span of our 2015 to 2022 data,
  but does not identify this particular dataset.
  \item \emph{Held out, TimesFM-3}: the fev-bench exclusion is stated specifically for the
  GIFT-Eval pretraining component. Because the model card does not specify its synthetic and
  augmented data further, TimesFM-3 has the least complete provenance documentation of the
  three models.
  \item \emph{Monash Traffic}: the pretraining corpora of Chronos-2 and TiRex contain other
  data from the same California PeMS traffic-sensor system, but not this dataset.
  \item \emph{Listed} establishes documented pretraining overlap, at the level of the dataset
  or, where marked \emph{source}, of its source. It does not imply that the specific windows
  used in our evaluation were presented during training.
\end{itemize}

\subsection{Windows and splits}
\label{app:splits}

\paragraph{Windows.}
Every example consists of a context of $C=512$ steps followed by a target of the next $H=64$ steps. Within each series, candidate forecast origins are spaced $H$ steps apart from the beginning of the series, so targets from different windows never overlap. A window is discarded if its context or target contains a missing value, or if its context is near-constant, with standard deviation below $10^{-6}$.

\paragraph{Temporal split.}
A series is eligible if it contains at least three valid windows. For each eligible series, the last valid window is reserved as its candidate test window, the second-to-last as its candidate validation window, and all earlier windows as candidate training windows. Thus, every training target precedes the validation target, which precedes the test target, with no target time step shared across splits. Each series selected for validation and test contributes exactly one window to each split.

\paragraph{Budgets.}
Each dataset uses 1,394 training windows and, by default, 262 validation and 262 test windows. Validation and test series are sampled from the eligible series; for eligible series not selected, the final two windows remain unused. Training windows are drawn from all retained eligible series and subsampled to 1,394 by round-robin over series, so every series contributes one window before any contributes a second and long series cannot dominate. Four datasets deviate from the default validation/test budget:
\begin{itemize}[leftmargin=*]
  \item SG Carpark and BOOM use all 354 eligible series, giving 354 validation and 354 test windows each.
  \item Coastal T-S uses all 48 eligible series, giving 48 validation and 48 test windows.
  \item SZ Taxi uses all 156 eligible series, giving 156 validation and 156 test windows.
\end{itemize}

\paragraph{Series cap.}
London Smart Meters, M5, and Wiki Daily have 5,555, 28,491, and 100,000 eligible series, respectively, more than the 1,394 training-window budget can cover. Before constructing the splits, we retain a fixed-seed random sample of 1,394 eligible series without using their data values. No series-level cap is applied to the remaining datasets.

\paragraph{Eligibility.}
Most datasets retain all series after filtering. The exceptions are Wind Farms, with 328 of 337 series eligible, and BOOM, with 354 of the 356 selected metric queries eligible.

\paragraph{Bootstrap clusters.}
Confidence intervals are obtained by resampling whole clusters. The cluster is one series by default; for SG Carpark it is one carpark, for BOOM one metric query, and for Coastal T-S one station, grouping its temperature and salinity series and yielding 24 independent clusters. Datasets with fewer independent clusters consequently have greater bootstrap uncertainty.

\paragraph{Identical windows across models.}
All three models are evaluated on exactly the same windows. For each dataset, the windows are fixed once and stored as a reference containing the split counts, the series identity of every test window, and a checksum of the test contexts. Every model run is checked against this reference and stops on any mismatch. All random choices use seed 0.

\subsection{Target spaces}
  \label{app:targets}

  Each model normalizes its input internally, and we train its readouts in the same normalized space
  as its native head. All statistics come from the context $x_{1:C}$ only, never from the future $y_{1:H}$.
  \begin{itemize}[leftmargin=*]
    \item \textbf{Chronos-2:} $z_t=\operatorname{arcsinh}\bigl((y_t-\mu)/\sigma\bigr)$, with $\mu$ and $\sigma$ the context mean and standard deviation.
    \item \textbf{TimesFM-3:} $z_t=(y_t-\tau_t-\mu_{15})/\sigma_{15}$, where $\tau_t$ is a least-squares linear trend extrapolated from the context (set to zero unless the residual standard deviation is below
  half the context standard deviation), and $\mu_{15},\sigma_{15}$ are the running statistics the model computes at token 15, whose state its head reads.
    \item \textbf{TiRex:} $z_t=(y_t-m_k)/s_k$ for the 32 steps forecast by pass $k\in\{1,2\}$, with $m_k,s_k$ the mean and standard deviation of that pass's context ($s_k=0$ is replaced by
  $|m_k|+10^{-5}$).
  \end{itemize}
  Predictions are mapped back to raw units by the models' own inverses: $\mu+\sigma\sinh(\hat z_t)$ for Chronos-2, $\sigma_{15}\hat z_t+\mu_{15}$ clipped to the model's value range and then $+\tau_t$ for
  TimesFM-3, and $s_k\hat z_t+m_k$ for TiRex. All three are monotone, so they preserve quantile order, and MASE is computed in raw units (Appendix~\ref{app:metrics}). Because the target spaces differ,
  readout losses are compared only within a model.

\subsection{Readout training}
\label{app:readout_training}

\paragraph{Inputs.}
At each depth, readout features are standardized per dimension using means and standard
deviations estimated from the training windows only. For Chronos-2 and TiRex, one
standardization is shared across the $K$ readout states, matching the shared readout.

\paragraph{Objective.}
Readouts are trained with pinball loss over the nine quantile levels and all $H=64$ horizon
steps in each model's target space (Appendix~\ref{app:targets}). We report mean pinball loss for
all three models. The Chronos-2 optimization code sums rather than averages over quantile
levels, differing from the reported convention only by a constant factor ($2Q$); this does not
affect weight-decay selection or the tunnel entrance.

\paragraph{Optimization.}
Each readout is a linear layer trained with full-batch AdamW for 300 epochs at learning rate
$10^{-2}$, with weight decay applied to the weights but not the bias. We retain the final-epoch
parameters without early stopping. For each model, dataset, and depth, every weight-decay
candidate starts from the same seeded initialization.

\paragraph{Weight-decay selection.}
For every model, dataset, and depth, we train one readout for each weight decay in
\[
\{10^{-5},10^{-4},10^{-3},10^{-2},10^{-1},0.3,1,3,10,30,45,65,90\}
\]
on the full training split and select the value with the lowest validation loss. The selected
readout is not retrained on training and validation data jointly, and the test split is used
only for final evaluation. We restrict the grid to values satisfying $\eta\lambda<1$, where
$\eta$ is the learning rate and $\lambda$ the weight decay: AdamW's decoupled decay multiplies
the current weights by $1-\eta\lambda$ at each update, which reaches zero at $\eta\lambda=1$ and
becomes negative beyond it.

\paragraph{Largest weight decay.}
When the largest value in the grid is selected, we compare its validation loss with the best
constant-forecast loss, obtained in closed form from the per-step training quantiles. This
provides a diagnostic for distinguishing an effectively uninformative readout from a
regularization optimum that may lie beyond the search range.

\paragraph{Exception: LOOP Seattle.}
For LOOP Seattle, the largest weight decay was selected in the $10^{-2}$ runs. We therefore
use learning rate $10^{-3}$ for all three models and augment the grid with
$\{150,300,450,650,900,950,975,990\}$, keeping $\eta\lambda<1$. All other readout-training
settings remain unchanged. As for every other dataset, the no-information reference is the
closed-form constant forecast.

\subsection{Alignment training}
\label{app:alignment}

\paragraph{Adapter and objective.} The adapter of Section~\ref{sec:adapter} has
$\Delta_\ell\in\mathbb{R}^{d\times d}$ and $b_\ell\in\mathbb{R}^{d}$, i.e.\ $d^2+d$ parameters per
depth: 0.59M for Chronos-2, 1.64M for TimesFM-3 and 0.26M for TiRex. Like the readout, one adapter
is shared across all states read by the head (the four forecast slots of Chronos-2 and the two pass
states of TiRex). It minimizes
\begin{equation}
\frac{1}{N_{\mathrm{train}}}\sum_{i=1}^{N_{\mathrm{train}}}
\bigl\| g\bigl(A_\ell(h_i^{(\ell)})\bigr) - g\bigl(h_i^{(L)}\bigr) \bigr\|_2^2
\label{eq:adapter_objective}
\end{equation}
over the training windows, with both forecasts in the model's normalized output space
(Appendix~\ref{app:targets}). The targets are the full model's own forecasts, so no future values
are used for fitting or selection. The selected adapter is not retrained, and the test split is read
once.

\paragraph{Chronos-2 and TiRex.} Their heads are nonlinear, so the adapter is fitted iteratively:
full-batch AdamW at learning rate $10^{-3}$ for 1,500 epochs, starting from $\Delta_\ell=0$,
$b_\ell=0$ (the hard cut). Weight decay acts on $\Delta_\ell$ only and therefore pulls the adapter
toward the identity. We train one adapter per weight decay in $\{0, 0.1, 1, 10, 30, 100\}$,
evaluate the validation error every 25 epochs, and keep the weight decay and epoch with the lowest
validation error. Epoch 0, the hard cut, is always a candidate, so the adapter is used only if it
improves on the hard cut on validation data. These optimizer settings were fixed before the full
runs, on validation data only, at the tunnel entrances of Chronos-2 and TiRex on Electricity: of
three settings (learning rate $10^{-2}$ for 300 epochs, $3\times10^{-3}$ for 1,000, and $10^{-3}$
for 1,500), only the last converged for every fit, i.e.\ selected a checkpoint before 90\% of the
epochs and a weight decay below the grid maximum.

\paragraph{TimesFM-3.} Its head is a single linear map $W$, so the objective reduces to a ridge
problem in the metric $W^\top W$ with a closed-form solution:
\begin{equation}
\min_{\Delta_\ell,\,b_\ell}\;\sum_{i}\bigl\|(I+\Delta_\ell)h_i^{(\ell)}+b_\ell-h_i^{(L)}\bigr\|_{W^\top W}^2
+\lambda\,\|I+\Delta_\ell\|_F^2,
\label{eq:adapter_ridge}
\end{equation}
with $b_\ell$ unpenalized and $\lambda=\kappa\,g_{\max}m_{\max}$, where $g_{\max}$ and $m_{\max}$
are the largest eigenvalues of the centred training Gram matrix and of $W^\top W$, and
$\kappa\in\{10^{-12},10^{-11},\dots,10^{2}\}$. The penalty shrinks the linear map
$A_\ell=I+\Delta_\ell$ toward zero; its limit, a constant forecast from the mean full-model state, is
an explicit candidate. The candidate with the lowest validation error in the same metric is kept,
with ties going to the larger $\lambda$. The submitted version of this paper penalized
$\lambda\|\Delta_\ell\|_F^2$ instead, shrinking toward the hard cut, with the hard cut as the
explicit candidate; for TimesFM-3 this overfits, and Appendix~\ref{sec:app_timesfm3_ridge} compares
the two.

\paragraph{Nonlinear variant.} The nonlinear variant adds a bottleneck branch,
$A_\ell(h)=h+\Delta_\ell h+b_\ell+W_2\,\mathrm{GELU}(W_1\,\mathrm{LN}(h)+c)$, where $\mathrm{LN}$ is
a layer normalization without learned parameters, $W_1\in\mathbb{R}^{64\times d}$ is initialized
randomly with a fixed seed and $W_2\in\mathbb{R}^{d\times 64}$ at zero, so the branch contributes
nothing at initialization. It uses the same objective and iterative protocol for all three models,
with weight decay also on $W_1$ and $W_2$. We tested it on Electricity at 25, 50 and 75\% of depth
under an adoption rule fixed in advance and evaluated on validation data only; the rule was not met,
and the full run on all 14 datasets agrees (Appendix~\ref{app:compatibility}). For TimesFM-3, the
iterative fit, with or without the branch, does not fit even the training windows (on Electricity,
training $R^2$ between $-314$ and $-6$ against the full model's forecast). The closed-form affine
adapter anchored at the hard cut fits them almost exactly ($R^2\approx0.997$) but not the validation
windows ($R^2$ between $-3.0$ and $-0.6$), whereas the zero-anchored closed form of
Equation~\ref{eq:adapter_ridge} reaches training $R^2$ between $0.96$ and $0.99$ and validation $R^2$
between $0.51$ and $0.82$. We therefore use the zero-anchored closed-form adapter for all TimesFM-3
results.


\subsection{Metrics and uncertainty}
\label{app:metrics}

\paragraph{Pinball loss.} For quantile level $\tau$, the pinball loss of a prediction $\hat z$
for a target $z$ is $\rho_\tau(z-\hat z)$ with $\rho_\tau(u)=\max\{\tau u,(\tau-1)u\}$. The
readout loss of a window is the mean of $\rho_\tau$ over the nine quantile levels and the $H$
horizon steps, in the model's target space (Appendix~\ref{app:targets}); a dataset's loss is the
mean over its windows. This loss selects the weight decay and defines the forecasting tunnel.

\paragraph{MASE.} Forecast accuracy in raw units is measured by the mean absolute scaled error
\citep{HYNDMAN2006679} of the median forecast. For a window with context $x$, future $y$ and
median forecast $\hat y$:
\begin{equation}
\mathrm{MASE}=\frac{\frac{1}{H}\sum_{t=1}^{H}\lvert y_t-\hat y_t\rvert}
{\frac{1}{C-m}\sum_{t=m+1}^{C}\lvert x_t-x_{t-m}\rvert},
\label{eq:mase}
\end{equation}
where $m$ is the dataset's seasonal period (Table~\ref{tab:datasets}). The denominator, the error
of a seasonal naive forecast, is computed from the window's own context, so every window has its
own scale and no future value enters it. It is floored at $10^{-8}$. A dataset's MASE is the mean
over its test windows. The full model, every readout and every truncated model are scored on the
same windows with the same denominators, so their MASE values are directly comparable. This
in-context scaling differs from benchmarks such as GIFT-Eval and fev-bench, which scale by the
series' history before the forecast, so our absolute MASE values are not comparable to published
leaderboards.

\paragraph{Confidence intervals.} Windows of the same series are correlated, so we resample whole
clusters: series, or carparks, metric queries and stations for SG Carpark, BOOM and Coastal
T-S (Appendix~\ref{app:splits}). Each of $B=5{,}000$ replicates draws as many clusters as the
dataset has, with replacement, and recomputes the window mean. The 95\% interval spans the 2.5th
to 97.5th percentiles of the replicates. The resampled clusters depend only on the dataset and
the seed (0), so every depth, model pathway and metric of a dataset uses the same replicates.
Comparisons are therefore paired: the interval of a difference, such as a depth against the final
block, or of a ratio, such as a truncated model's MASE over the full model's, is computed within
each replicate.

\subsection{Latency protocol}
\label{app:latency}

\paragraph{What is timed.}
We time each model's public forecasting interface end to end, from a context array in host memory
to quantile forecasts returned to host memory: \texttt{predict\_quantiles} for Chronos-2,
\texttt{predict\_batch} for TimesFM-3, and \texttt{forecast} with default two-pass inference for
TiRex. This is our headline measurement. As a secondary measurement, we time the forward pass alone
with inputs already on the GPU. Each configuration is measured at batch sizes 1 (latency) and 256
(throughput setting). We time the full model and, at every depth, the hard cut, the aligned model
and the readout used as a replacement head. Execution time depends on an adapter's shape rather
than its learned values, so adapters are timed with random weights.

\paragraph{Inputs.}
For complete contexts of fixed length, the timed computation does not depend on the input values,
so we use synthetic random-walk contexts of length $C=512$ rather than the 14 evaluation datasets.
Four fixed batches are cycled so that consecutive calls receive different inputs.

\paragraph{Procedure.}
All timings use one full NVIDIA A100 GPU rather than a partitioned device, with TF32 disabled and
no compilation or mixed-precision inference. Model parameters and surrounding computation use
float32; TiRex's recurrence internally uses bfloat16, as in all our experiments. Each configuration
runs in a fresh process, so removed blocks leave no cached state, and configurations run in random
order. After 20 warm-up calls, we time 100 calls with GPU synchronization before and after each
call and record the median. Only truncated configurations that passed our verification of physical
truncation are timed.

\paragraph{Speedup.}
Within each job, speedup is the full model's median runtime divided by the configuration's median
runtime. The full model is timed at the beginning, middle and end of each job, and the median of the
three is the baseline; the three differ by a median of 1.0\% and at most 5.4\%. We repeat the
procedure in three independent jobs (TiRex's depths are split over two jobs per repetition, each
with its own baseline) and report the median of the per-job speedups.

\paragraph{Full-model runtimes.} At batch size 1, one end-to-end forecast takes 19.4\,ms for
Chronos-2, 97.5\,ms for TimesFM-3 and 499\,ms for TiRex. For TimesFM-3, the forward pass alone takes
49.3\,ms, so about half of its end-to-end time is spent outside the backbone and cannot be removed
by truncation; this limits its end-to-end speedup at batch size 1.

\paragraph{Physical versus offline truncation.}
Forecasting performance along the frontier (Figure~\ref{fig:h4_frontier} and
Table~\ref{tab:h4_frontier_full}) is computed by applying the aligned pathway to the cached
representations of each depth, while speed is measured on physically truncated models. We verify
that both describe the same model. At every depth of every model, the physically truncated network
reproduces the full model's state at that depth bitwise, on synthetic inputs and on 256 Electricity
windows against the cached features; the removed blocks are never executed; and the physically
truncated aligned model reproduces the offline aligned forecast within a relative tolerance of
$10^{-4}$. In addition, 419 truncated configurations, covering every model-dataset pair, were
evaluated physically on the complete test splits; their test MASE agrees with the offline pathway
to within $7.9\times10^{-5}$ relative, so no reported result depends on which pipeline is used.
Individual forecast values agree less tightly. For TiRex, even the unmodified model differs between
the two pipelines by up to 11\% of the mean forecast magnitude on single values, consistent with its
bfloat16 recurrence; for TimesFM-3, the aligned pathway differs by up to 5\% while the unmodified
model agrees. The element-wise check we had specified in advance for these test-split evaluations
(relative tolerance $10^{-4}$) therefore rejected 96 of the 419 configurations. We replaced it,
uniformly for all models, by the requirement that test MASE agree within $10^{-4}$ relative, which
every configuration satisfies. These checks cover the adapters of the submitted version; the
zero-anchored TimesFM-3 adapters have the same shape and therefore the same latency, and their
physical check is reported in Appendix~\ref{sec:app_timesfm3_ridge}.

\subsection{Compute}
\label{app:compute}

All experiments ran on a shared academic GPU cluster. Each GPU job used one NVIDIA A100 (40\,GB);
building windows and computing representation geometry need no GPU. The software stack was
Python~3.11, PyTorch~2.12 with CUDA~13.2, \texttt{chronos-forecasting}~2.3.1, \texttt{timesfm}~3.0.2
and \texttt{tirex-ts}~1.4.2. The pretrained backbones are never trained, so all computation is
inference plus the fitting of small readouts and adapters. Table~\ref{tab:compute} reports GPU time
by stage. Over half of it is latency measurement; the probing analyses themselves take under
3~GPU-hours. The nonlinear adapter, which we did not adopt, accounts for 18.4 of the 48.9
GPU-hours. The experiments added during the review period used a further 20.2 A100-hours for
Appendices~\ref{sec:app_content_mean}, \ref{sec:app_lodo}, \ref{sec:app_source_diversity},
and~\ref{sec:app_noncontig}, not counting duplicate runs on another cluster whose results we do not
use; the closed-form TimesFM-3 fits of Appendix~\ref{sec:app_timesfm3_ridge} and their physical check
ran on CPUs.

\begin{table}[h]
  \centering
  \small
  \vspace{4pt}
  \caption{\textbf{Compute by stage}, in A100 GPU-hours from the scheduler's accounting, summed
  over all GPU jobs including failed and cancelled ones. Feature extraction runs inside the readout
  jobs. Alignment, truncation and latency include both the affine and the nonlinear adapter.}
  \label{tab:compute}
  \begin{tabular}{@{}lr@{}}
    \toprule
    Stage & GPU-hours \\
    \midrule
    Feature extraction, layerwise readouts ($Q=9$ and $Q=1$) and geometry & 2.9 \\
    Alignment (all depths, models and datasets) & 13.8 \\
    Truncation checks and evaluation & 0.5 \\
    Latency (three repetitions) & 26.0 \\
    Development, smoke tests and superseded runs & 5.7 \\
    \midrule
    Total & 48.9 \\
    \bottomrule
  \end{tabular}
\end{table}




\section{Additional Results}
\label{app:additional_results}

\subsection{Additional forecast-recoverability results}
\label{sec:app_recoverability}

\subsubsection{Complete layerwise forecasting curves}
\label{sec:app_recov_curves}

Figures~\ref{fig:app_chronos_q9}, \ref{fig:app_timesfm_q9} and~\ref{fig:app_tirex_q9} show the complete nine-quantile ($Q=9$) layerwise forecasting results for all three models and all 14 datasets. Tunnel entrances are selected from validation loss using the sustained 5\% criterion, while test loss is shown only for held-out evaluation.

\paragraph{Pretraining exposure and the tunnel entrance.}
  We find no consistent relationship between a dataset's documented pretraining exposure (Table~\ref{tab:provenance}; cells marked \emph{source} count as listed)  and
  where its tunnel begins. Within each model, the median normalized tunnel entrance
  $\ell_{\mathrm{tunnel}}/L$ for listed versus held-out or not-listed datasets is $0.33$ versus
  $0.25$ for Chronos-2, $0.05$ versus $0.25$ for TimesFM-3, and $0.83$ versus $0.67$ for TiRex. The
  direction therefore differs across models, and exact permutation tests do not distinguish the
  groups ($p\geq0.58$). 

\begin{figure}[p]
  \centering
  \includegraphics[width=\textwidth]{figures/appendix_chronos2_q9.pdf}
  \caption{\textbf{Chronos-2 layerwise Q=9 forecasting loss across all datasets.} The tunnel entrance is selected from validation loss using the sustained 5\% criterion. Bands show 95\% series-level cluster-bootstrap intervals ($B=5000$).}
  \label{fig:app_chronos_q9}
\end{figure}

\begin{figure}[p]
  \centering
  \includegraphics[width=\textwidth]{figures/appendix_timesfm3_q9.pdf}
  \caption{\textbf{TimesFM-3 layerwise Q=9 forecasting loss across all datasets.} The tunnel entrance is selected from validation loss using the sustained 5\% criterion. Bands show 95\% series-level cluster-bootstrap intervals ($B=5000$).}
  \label{fig:app_timesfm_q9}
\end{figure}

\begin{figure}[p]
  \centering
  \includegraphics[width=\textwidth]{figures/appendix_tirex_q9.pdf}
  \caption{\textbf{TiRex layerwise Q=9 forecasting loss across all datasets.} The tunnel entrance is selected from validation loss using the sustained 5\% criterion. Bands show 95\% series-level cluster-bootstrap intervals ($B=5000$).}
  \label{fig:app_tirex_q9}
\end{figure}


\subsubsection{Held-out performance at the tunnel entrance}
\label{sec:app_heldout_recoverability}

Table~\ref{tab:heldout_q9} reports held-out performance at each validation-selected tunnel entrance. Across all 42 model-dataset pairs, the median increase in $Q=9$ test loss relative to the final-block readout is 1.3\% (interquartile range $-0.7$\% to 4.3\%), and 34 of the 42 pairs stay within 5\% on test, both at the entrance and at every later block. By model, the median increase is 2.8\% for Chronos-2, 0.8\% for TimesFM-3 and 2.1\% for TiRex. The largest increase is 9.8\%, for TimesFM-3 on Uber TLC; negative values, where the entrance readout is better on test than the final-block readout, occur for all three models. The median-forecast MASE changes by a median of $+1.3$\%.

\begin{table}[p]
  \centering
  \caption{\textbf{Held-out performance at the validation-selected 5\% tunnel entrance.} $\Delta$ test gives the relative change in $Q=9$ test loss compared with the final-block readout, with a paired 95\% series-level cluster-bootstrap interval ($B=5000$). Negative values indicate lower test loss at the tunnel entrance. Entrances are selected using validation loss only; where the entrance is the final block, $\Delta$ test is zero by construction.}
  \vspace{4pt}
  \label{tab:heldout_q9}
  \scriptsize
  \setlength{\tabcolsep}{8pt}
  \begin{tabular}{@{}lccc@{\qquad}lccc@{}}
  \toprule
  Dataset & Entr. & $\Delta$ test (\%) & 95\% CI &
  Dataset & Entr. & $\Delta$ test (\%) & 95\% CI \\
  \midrule
  \multicolumn{8}{l}{\textbf{Chronos-2}} \\
  M4 & L8 & $+7.1$ & $[5.0, 9.3]$ &
  LOOP Seattle & L9 & $-6.1$ & $[-9.0, -3.1]$ \\
  Electricity & L9 & $+3.0$ & $[1.5, 4.7]$ &
  Electricity (15-minute) & L10 & $-1.0$ & $[-2.7, 0.7]$ \\
  Uber TLC & L4 & $+6.4$ & $[5.3, 7.5]$ &
  SZ Taxi & L2 & $+3.7$ & $[1.5, 6.6]$ \\
  Wind Farms & L3 & $+1.0$ & $[-3.9, 6.3]$ &
  London Smart Meters & L3 & $+3.3$ & $[1.7, 5.0]$ \\
  SG Carpark & L10 & $-4.4$ & $[-6.0, -2.8]$ &
  M5 & L0 & $+6.9$ & $[4.8, 9.2]$ \\
  Coastal T-S & L1 & $+2.5$ & $[-4.4, 10.0]$ &
  Wiki Daily & L2 & $+6.7$ & $[4.3, 9.4]$ \\
  BOOM & L3 & $+1.3$ & $[-1.9, 4.3]$ &
  Monash Traffic & L11 & $-0.2$ & $[-0.9, 0.6]$ \\
  \midrule
  \multicolumn{8}{l}{\textbf{TimesFM-3}} \\
  M4 & L16 & $-0.8$ & $[-2.8, 1.4]$ &
  LOOP Seattle & L14 & $-8.3$ & $[-11.5, -5.0]$ \\
  Electricity & L15 & $-1.2$ & $[-2.9, 0.6]$ &
  Electricity (15-minute) & L1 & $+1.1$ & $[-2.9, 5.3]$ \\
  Uber TLC & L1 & $+9.8$ & $[8.6, 11.1]$ &
  SZ Taxi & L1 & $+1.3$ & $[0.1, 2.6]$ \\
  Wind Farms & L0 & $+3.7$ & $[1.4, 14.8]$ &
  London Smart Meters & L1 & $+0.4$ & $[-0.9, 2.0]$ \\
  SG Carpark & L12 & $-3.1$ & $[-5.0, -1.3]$ &
  M5 & L0 & $+2.9$ & $[2.0, 4.0]$ \\
  Coastal T-S & L5 & $-6.3$ & $[-13.8, 2.0]$ &
  Wiki Daily & L0 & $+6.2$ & $[3.9, 9.0]$ \\
  BOOM & L0 & $-4.4$ & $[-9.2, -0.1]$ &
  Monash Traffic & L12 & $+1.9$ & $[-0.9, 4.9]$ \\
  \midrule
  \multicolumn{8}{l}{\textbf{TiRex}} \\
  M4 & L12 & $0.0$ & $[0.0, 0.0]$ &
  LOOP Seattle & L9 & $-1.1$ & $[-3.5, 1.5]$ \\
  Electricity & L11 & $+4.5$ & $[2.9, 6.0]$ &
  Electricity (15-minute) & L10 & $+2.8$ & $[1.5, 4.1]$ \\
  Uber TLC & L10 & $+1.4$ & $[0.9, 1.9]$ &
  SZ Taxi & L0 & $+4.8$ & $[3.1, 6.6]$ \\
  Wind Farms & L11 & $+0.7$ & $[0.2, 3.6]$ &
  London Smart Meters & L8 & $+5.2$ & $[4.1, 6.3]$ \\
  SG Carpark & L11 & $-1.8$ & $[-3.5, -0.2]$ &
  M5 & L0 & $+3.9$ & $[3.0, 4.8]$ \\
  Coastal T-S & L1 & $-0.7$ & $[-8.0, 7.0]$ &
  Wiki Daily & L1 & $+5.5$ & $[3.5, 7.7]$ \\
  BOOM & L8 & $+4.7$ & $[1.6, 7.7]$ &
  Monash Traffic & L12 & $0.0$ & $[0.0, 0.0]$ \\
  \bottomrule
  \end{tabular}
\end{table}

\subsubsection{Sensitivity to readout capacity}
\label{sec:app_q1_robustness}

We repeat the analysis with a median-only ($Q=1$) readout, whose output is a ninth of the size of the $Q=9$ readout, trained with the same settings as the $Q=9$ readout of each dataset (for LOOP Seattle, the lower learning rate and extended grid of Appendix~\ref{app:readout_training}). Table~\ref{tab:q1_q9_comparison} compares the tunnel entrances. Of the 42 pairs, 31 have exactly the same entrance, 37 differ by at most one block, and the median absolute difference is zero blocks. The largest difference is TimesFM-3 on Coastal T-S, where the $Q=1$ entrance moves from L5 to L18; Coastal T-S has the fewest bootstrap clusters (24). The next largest are TimesFM-3 on LOOP Seattle (four blocks earlier with $Q=1$) and Chronos-2 on M4 (three blocks earlier).

\begin{table}[t]
  \centering
  \caption{\textbf{$Q=1$ versus $Q=9$ sustained 5\% tunnel entrances.} Each cell gives the difference in entrance depth, $\Delta\ell=\ell_{Q=1}-\ell_{Q=9}$, with the $Q=1$ entrance in parentheses.}
  \vspace{4pt}
  \label{tab:q1_q9_comparison}
  \small
  \setlength{\tabcolsep}{3.5pt}
  \begin{tabular}{@{}lccc@{\qquad}lccc@{}}
  \toprule
  Dataset & Chronos-2 & TFM-3 & TiRex &
  Dataset & Chronos-2 & TFM-3 & TiRex \\
  \midrule
  M4 & $-3$ (L5) & $-2$ (L14) & 0 (L12) &
  LOOP Seattle & 0 (L9) & $-4$ (L10) & 0 (L9) \\
  Electricity & $+1$ (L10) & 0 (L15) & 0 (L11) &
  Electricity (15-minute) & 0 (L10) & 0 (L1) & 0 (L10) \\
  Uber TLC & $+1$ (L5) & 0 (L1) & 0 (L10) &
  SZ Taxi & 0 (L2) & 0 (L1) & 0 (L0) \\
  Wind Farms & 0 (L3) & 0 (L0) & 0 (L11) &
  London Smart Meters & $+1$ (L4) & 0 (L1) & 0 (L8) \\
  SG Carpark & $+2$ (L12) & $-1$ (L11) & 0 (L11) &
  M5 & 0 (L0) & 0 (L0) & 0 (L0) \\
  Coastal T-S & 0 (L1) & $+13$ (L18) & 0 (L1) &
  Wiki Daily & 0 (L2) & 0 (L0) & 0 (L1) \\
  BOOM & $-1$ (L2) & 0 (L0) & $+1$ (L9) &
  Monash Traffic & 0 (L11) & 0 (L12) & 0 (L12) \\
  \midrule
  \textbf{Exact match} & \textbf{8/14} & \textbf{10/14} & \textbf{13/14} &
  \textbf{$|\Delta\ell|\leq1$} & \textbf{12/14} & \textbf{11/14} & \textbf{14/14} \\
  \bottomrule
  \end{tabular}
\end{table}

\subsubsection{Sensitivity to the tunnel tolerance}
\label{sec:app_tunnel_sensitivity}

We repeat the analysis with $\epsilon\in\{0.02,0.05,0.08\}$, using the same validation curves. Table~\ref{tab:tunnel_sensitivity} lists the resulting entrances. The tunnel begins before the final block in 36 of the 42 pairs at $\epsilon=0.02$, 40 at $\epsilon=0.05$, and all 42 at $\epsilon=0.08$. Tightening the tolerance to 2\% moves the median entrance from 29\% to 67\% of model depth for Chronos-2 and from 5\% to 65\% for TimesFM-3, and loosening it to 8\% moves TiRex from 79\% to 67\%.

\begin{table}[p]
  \centering
  \caption{\textbf{Sensitivity of the sustained tunnel entrance to $\epsilon$.} Each entry gives the selected entrance for $\epsilon=0.02$, $0.05$, and $0.08$. Median depth is the median of $\ell/L$ over the 14 datasets.}
  \vspace{4pt}
  \label{tab:tunnel_sensitivity}
  \small
  \setlength{\tabcolsep}{3.5pt}
  \begin{tabular}{@{}lccc@{\qquad}lccc@{}}
  \toprule
  Dataset & 2\% & 5\% & 8\% &
  Dataset & 2\% & 5\% & 8\% \\
  \midrule
  \multicolumn{8}{l}{\textbf{Chronos-2}} \\
  M4 & L11 & L8 & L5 & LOOP Seattle & L12 & L9 & L8 \\
  Electricity & L11 & L9 & L5 & Electricity (15-minute) & L11 & L10 & L9 \\
  Uber TLC & L5 & L4 & L3 & SZ Taxi & L4 & L2 & L1 \\
  Wind Farms & L11 & L3 & L3 & London Smart Meters & L5 & L3 & L3 \\
  SG Carpark & L11 & L10 & L10 & M5 & L1 & L0 & L0 \\
  Coastal T-S & L3 & L1 & L1 & Wiki Daily & L3 & L2 & L1 \\
  BOOM & L4 & L3 & L2 & Monash Traffic & L12 & L11 & L11 \\
  \textbf{Median depth} & \textbf{67\%} & \textbf{29\%} & \textbf{25\%} &
  \textbf{Before final} & \textbf{12/14} & \textbf{14/14} & \textbf{14/14} \\
  \midrule
  \multicolumn{8}{l}{\textbf{TimesFM-3}} \\
  M4 & L20 & L16 & L13 & LOOP Seattle & L17 & L14 & L1 \\
  Electricity & L16 & L15 & L15 & Electricity (15-minute) & L11 & L1 & L1 \\
  Uber TLC & L13 & L1 & L1 & SZ Taxi & L5 & L1 & L0 \\
  Wind Farms & L15 & L0 & L0 & London Smart Meters & L1 & L1 & L1 \\
  SG Carpark & L13 & L12 & L4 & M5 & L1 & L0 & L0 \\
  Coastal T-S & L18 & L5 & L4 & Wiki Daily & L2 & L0 & L0 \\
  BOOM & L0 & L0 & L0 & Monash Traffic & L15 & L12 & L11 \\
  \textbf{Median depth} & \textbf{65\%} & \textbf{5\%} & \textbf{5\%} &
  \textbf{Before final} & \textbf{13/14} & \textbf{14/14} & \textbf{14/14} \\
  \midrule
  \multicolumn{8}{l}{\textbf{TiRex}} \\
  M4 & L12 & L12 & L11 & LOOP Seattle & L10 & L9 & L8 \\
  Electricity & L11 & L11 & L11 & Electricity (15-minute) & L10 & L10 & L9 \\
  Uber TLC & L11 & L10 & L8 & SZ Taxi & L7 & L0 & L0 \\
  Wind Farms & L11 & L11 & L8 & London Smart Meters & L11 & L8 & L2 \\
  SG Carpark & L12 & L11 & L10 & M5 & L1 & L0 & L0 \\
  Coastal T-S & L4 & L1 & L1 & Wiki Daily & L8 & L1 & L0 \\
  BOOM & L11 & L8 & L1 & Monash Traffic & L12 & L12 & L11 \\
  \textbf{Median depth} & \textbf{92\%} & \textbf{79\%} & \textbf{67\%} &
  \textbf{Before final} & \textbf{11/14} & \textbf{12/14} & \textbf{14/14} \\
  \bottomrule
  \end{tabular}
\end{table}


% Appendix B.1.5: architecture-aware versus generic content-mean readouts (Phase 2d).
% Moved verbatim from main.tex during the rebuttal revision; \input in Appendix B.1.

\subsubsection{Architecture-aware versus generic content-mean readouts}
\label{sec:app_content_mean}

In Section~\ref{sec:layerwise_readouts}, we probe the hidden states used by each model's native forecasting head. To investigate how this choice affects our results, we repeat the analysis with a more generic content-mean readout, which averages the hidden states of all context tokens at each depth and applies a linear map $\mathbb{R}^{d}\to\mathbb{R}^{Q\times H}$ to predict the full forecasting horizon. Unlike our architecture-aware readouts, this procedure does not depend on which states the model uses to produce its forecasts. For Chronos-2, we average the 32 context-patch tokens, excluding the forecast slots and register token; for TimesFM-3, the 16 context tokens; and for TiRex, the 16 context tokens of the first forecasting pass, excluding padding. All readouts use the same windows, targets, and training procedure described in Appendix~\ref{app:readout_training}.

The generic readout also differs from our original readout in how the selected states are mapped to forecasts. To separate this effect from the choice of hidden states, we additionally consider a \emph{native-state mean} readout, which averages the native forecast-producing states and applies the same linear map as the content-mean readout. The two mean readouts therefore have the same architecture and number of parameters, differing only in which hidden states are averaged. For TimesFM-3, the readout architecture is already identical, so no additional control is needed.

Table~\ref{tab:content_mean_loss} and Figure~\ref{fig:app_content_mean_curves} compare the forecasting performance of these readouts across all 14 datasets. At the final block, content-mean pooling increases the median test loss by 27.2\%, 15.9\%, and 20.9\% for Chronos-2, TimesFM-3, and TiRex, respectively, relative to our architecture-aware readouts. The paired 95\% confidence intervals exclude zero in the direction of higher loss on 12, 14, and 13 of the 14 datasets. Even at its best validation-selected depth, the content-mean readout remains 21.6\%, 10.7\%, and 19.2\% worse than the architecture-aware final-block readout. Comparing readouts with identical architectures and parameter counts, changing only the selected states increases median final-block loss by 11.1\%, 15.9\%, and 15.6\%, respectively. Thus, the choice of which hidden states to probe contributes substantially to the difference in forecasting performance.

We next examine whether generic pooling changes the location of the forecasting tunnel. Table~\ref{tab:content_mean_entrances} compares the sustained 5\% tunnel entrances for both readouts. The content-mean tunnel generally begins earlier: the median relative entrance depth moves from 29\% to 8\% for Chronos-2 and from 79\% to 21\% for TiRex, while remaining at 5\% for TimesFM-3. However, these earlier entrances do not necessarily indicate earlier recovery of accurate forecasts. Each tunnel is defined relative to its own final-block readout, and content-mean pooling produces substantially higher forecasting losses. In fact, its final-block test loss exceeds that of the architecture-aware readout already at the latter's tunnel entrance on 13, 14, and 13 of the 14 datasets, respectively. The earlier entrances therefore largely reflect the saturation of a weaker readout.

Overall, these results show that the choice of hidden states can substantially affect both forecasting performance and the apparent location of the tunnel. Simply averaging intermediate representations can suggest earlier saturation, even when the resulting forecasts remain considerably less accurate. Nevertheless, the qualitative finding that the tunnel begins before the final block remains unchanged, holding for all 42 model-dataset pairs under content-mean pooling, compared with 40/42 using our architecture-aware readouts.


\begin{table}[t]
\centering
\caption{\textbf{Final-block forecasting loss for architecture-aware and mean-pooled readouts.}
Native: absolute $Q=9$ test loss in each model's normalized target space. All dataset-level percentage changes are relative to the original architecture-aware (Native) readout. The final row instead compares content mean with native-state mean, holding the readout architecture and parameter count fixed; for TimesFM-3, native and native-state mean are identical. Medians are computed across 14 datasets from the corresponding per-dataset relative changes. $^{*}$ indicates that the paired 95\% cluster-bootstrap confidence interval excludes zero for the comparison against Native.}
\vspace{4pt}
\label{tab:content_mean_loss}
\scriptsize
\setlength{\tabcolsep}{3pt}
\resizebox{\textwidth}{!}{%
\begin{tabular}{@{}lccc cc ccc@{}}
\toprule
& \multicolumn{3}{c}{Chronos-2}
& \multicolumn{2}{c}{TimesFM-3}
& \multicolumn{3}{c}{TiRex} \\
\cmidrule(lr){2-4}
\cmidrule(lr){5-6}
\cmidrule(l){7-9}
Dataset
& Native & Nat.\ mean & Content mean
& Native & Content mean
& Native & Nat.\ mean & Content mean \\
\midrule
M4
& 0.099 & $+18.6^{*}$ & $+46.1^{*}$
& 0.132 & $+27.4^{*}$
& 0.106 & $+14.0^{*}$ & $+58.5^{*}$ \\
Electricity
& 0.118 & $+25.4^{*}$ & $+59.5^{*}$
& 0.169 & $+19.8^{*}$
& 0.151 & $+14.5^{*}$ & $+57.8^{*}$ \\
Uber TLC
& 0.179 & $+24.9^{*}$ & $+42.8^{*}$
& 0.263 & $+21.3^{*}$
& 0.248 & $+12.9^{*}$ & $+28.6^{*}$ \\
Wind Farms
& 0.159 & $+16.2^{*}$ & $+25.4^{*}$
& 0.608 & $+7.5^{*}$
& 0.589 & $+1.8^{*}$ & $+9.3^{*}$ \\
SG Carpark
& 0.101 & $+0.6$ & $+51.8^{*}$
& 0.134 & $+41.4^{*}$
& 0.121 & $+4.5^{*}$ & $+64.2^{*}$ \\
Coastal T-S
& 0.199 & $+3.8^{*}$ & $+5.7$
& 0.326 & $+15.9^{*}$
& 0.293 & $+0.8$ & $+11.5$ \\
BOOM
& 0.131 & $+6.0^{*}$ & $+10.8^{*}$
& 0.238 & $+13.3^{*}$
& 0.212 & $+1.3^{*}$ & $+17.0^{*}$ \\
LOOP Seattle
& 0.191 & $-2.6^{*}$ & $-5.7^{*}$
& 0.234 & $+15.9^{*}$
& 0.239 & $+0.4$ & $+20.9^{*}$ \\
Electricity (15-min)
& 0.165 & $+14.2^{*}$ & $+33.6^{*}$
& 0.204 & $+57.2^{*}$
& 0.201 & $+8.0^{*}$ & $+38.5^{*}$ \\
SZ Taxi
& 0.237 & $+3.8^{*}$ & $+6.8^{*}$
& 0.323 & $+4.7^{*}$
& 0.313 & $+2.0^{*}$ & $+6.3^{*}$ \\
London Smart Meters
& 0.148 & $+11.4^{*}$ & $+29.1^{*}$
& 0.216 & $+12.6^{*}$
& 0.203 & $+4.5^{*}$ & $+20.9^{*}$ \\
M5
& 0.220 & $+2.0^{*}$ & $+4.4^{*}$
& 0.293 & $+5.8^{*}$
& 0.288 & $+1.1^{*}$ & $+3.2^{*}$ \\
Wiki Daily
& 0.183 & $+4.4^{*}$ & $+10.4^{*}$
& 0.247 & $+8.3^{*}$
& 0.242 & $+1.7^{*}$ & $+8.5^{*}$ \\
Monash Traffic
& 0.084 & $+15.4^{*}$ & $+53.8^{*}$
& 0.119 & $+31.5^{*}$
& 0.099 & $+13.7^{*}$ & $+57.0^{*}$ \\
\midrule
\textbf{Median change vs.\ Native}
& -- & $\mathbf{+8.7\%}$ & $\mathbf{+27.2\%}$
& -- & $\mathbf{+15.9\%}$
& -- & $\mathbf{+3.2\%}$ & $\mathbf{+20.9\%}$ \\
\textbf{Content vs.\ Nat.\ mean}
& -- & -- & $\mathbf{+11.1\%}$
& -- & $\mathbf{+15.9\%}$
& -- & -- & $\mathbf{+15.6\%}$ \\
\bottomrule
\end{tabular}%
}
\end{table}


\begin{table}[t]
\centering
\caption{\textbf{Sustained 5\% forecasting-tunnel entrances for architecture-aware and content-mean readouts.}
Each pair gives the entrance of the architecture-aware (Native) and generic content-mean (Pooled) readout, selected using validation loss relative to that readout's own final block. L0 denotes the representation before the first block. Median depth is the median entrance as a fraction of total model depth.}
\vspace{4pt}
\label{tab:content_mean_entrances}
\small
\setlength{\tabcolsep}{4pt}
\begin{tabular}{@{}lcc cc cc@{}}
\toprule
& \multicolumn{2}{c}{Chronos-2}
& \multicolumn{2}{c}{TimesFM-3}
& \multicolumn{2}{c}{TiRex} \\
\cmidrule(lr){2-3}
\cmidrule(lr){4-5}
\cmidrule(l){6-7}
Dataset
& Native & Pooled
& Native & Pooled
& Native & Pooled \\
\midrule
M4 & L8 & L9 & L16 & L18 & L12 & L7 \\
Electricity & L9 & L10 & L15 & L1 & L11 & L10 \\
Uber TLC & L4 & L0 & L1 & L1 & L10 & L8 \\
Wind Farms & L3 & L2 & L0 & L1 & L11 & L1 \\
SG Carpark & L10 & L1 & L12 & L1 & L11 & L9 \\
Coastal T-S & L1 & L1 & L5 & L0 & L1 & L1 \\
BOOM & L3 & L2 & L0 & L1 & L8 & L5 \\
LOOP Seattle & L9 & L8 & L14 & L1 & L9 & L0 \\
Electricity (15-min) & L10 & L10 & L1 & L1 & L10 & L4 \\
SZ Taxi & L2 & L0 & L1 & L0 & L0 & L0 \\
London Smart Meters & L3 & L0 & L1 & L0 & L8 & L0 \\
M5 & L0 & L0 & L0 & L0 & L0 & L0 \\
Wiki Daily & L2 & L1 & L0 & L0 & L1 & L0 \\
Monash Traffic & L11 & L1 & L12 & L1 & L12 & L11 \\
\midrule
\textbf{Median depth}
& \textbf{29\%} & \textbf{8\%}
& \textbf{5\%} & \textbf{5\%}
& \textbf{79\%} & \textbf{21\%} \\
\textbf{Before final}
& \textbf{14/14} & \textbf{14/14}
& \textbf{14/14} & \textbf{14/14}
& \textbf{12/14} & \textbf{14/14} \\
\bottomrule
\end{tabular}
\end{table}

\begin{figure}[t]
\centering
\includegraphics[width=\textwidth]{figures/pooled_curves_main.pdf}
\caption{\textbf{Architecture-aware versus content-mean readouts for Electricity, Uber TLC, and London Smart Meters.}
Test $Q=9$ forecasting loss across depth for the architecture-aware readout, generic content-mean pooling, and the native-state mean control. Circles mark the sustained 5\% tunnel entrances, selected using validation loss relative to each readout's final block. Shaded bands show 95\% series-level cluster-bootstrap confidence intervals ($B=5000$).}
\label{fig:app_content_mean_curves}
\end{figure}


% Appendix B.1.6: the forecasting tunnel referenced to the full model (review period).
% Numbers: experiments/make_rebuttal_numbers.py, section A (results/rebuttal_numbers/).

\subsubsection{Tunnel referenced to the full model}
\label{sec:app_native_reference}

The tunnel of Equation~\ref{eq:tunnel_entrance} is referenced to the final-block readout, not to the full model. Here we replace the reference by the full model's validation loss $\mathcal{R}^{\mathrm{val}}_{\mathrm{full}}$, computed in the same target space, and keep the sustained criterion unchanged: $\mathcal{R}^{\mathrm{val}}_{j}\le(1+\epsilon)\,\mathcal{R}^{\mathrm{val}}_{\mathrm{full}}$ for every $j\in\{\ell,\dots,L\}$. Because the final-block readout is part of every suffix, such a tunnel exists only where the final-block readout is itself within $\epsilon$ of the full model.

Table~\ref{tab:native_reference} shows that this is rare. At the final block, the readout's validation loss exceeds the full model's by a median 17.7\%, 20.3\%, and 12.7\% for Chronos-2, TimesFM-3, and TiRex, and is within 5\% in only 6 of the 42 pairs: Chronos-2 and TiRex on SZ Taxi and Coastal T-S, TimesFM-3 on Coastal T-S, and TiRex on Wiki Daily. In all six, the tunnel referenced to the full model begins before the final block. At $\epsilon=0.10$, 15 pairs have such a tunnel, all beginning before the final block; at $\epsilon=0.02$, two pairs have one, and in only one of them does it begin before the final block. The tunnel of the main analysis therefore measures where later blocks stop improving a fixed linear readout; the gap between that readout and the full model is a property of the readout, present at every depth.

The aligned pathway, unlike a readout, reproduces the full model at the final block, so the same criterion applied to its validation loss asks from which depth the aligned truncated model stays within $\epsilon$ of the full model. Such a depth exists before the final block in 33 of 42 pairs at $\epsilon=0.05$ and 39 of 42 at $\epsilon=0.10$, using the zero-anchored TimesFM-3 adapter of Appendix~\ref{sec:app_timesfm3_ridge} (27 and 37 with the identity-anchored adapter of the submission). For Chronos-2 and TiRex, this depth is rank-correlated with the tunnel entrance across datasets (Spearman 0.81 and 0.80), and the two differ by at most one block on 9 and 11 of 14 datasets. For TimesFM-3 there is no such relation ($-0.02$): its aligned pathway stays within 5\% of the full model only from a median of L18.5, wherever the tunnel begins.

\paragraph{Entrances at L0.} Seven pairs enter the tunnel at L0. In four of them, M5 for all three models and TiRex on SZ Taxi, the final-block readout's validation loss is within 5\% of the best constant forecast (ratios 0.95 to 1.01; Appendix~\ref{app:readout_training}), so an L0 entrance indicates that no depth adds forecasting performance a linear readout can use. These pairs are marked in Table~\ref{tab:tunnel_q9}. In the other three, TimesFM-3 on BOOM, Wind Farms, and Wiki Daily, the final-block readout is clearly better than the constant forecast (ratios 0.72 to 0.88) and the L0 readout matches it.

\begin{table}[t]
\centering
\caption{\textbf{Tunnels referenced to the full model.} Validation $Q=9$ loss. Final block: readout loss relative to the full model (median over 14 datasets) and the number of datasets within 5\%. Readout and aligned pathway: number of datasets whose sustained tunnel, referenced to the full model's validation loss, begins before the final block, at $\epsilon=0.02$, $0.05$, and $0.10$. Spearman: rank correlation across datasets between the readout tunnel entrance of Table~\ref{tab:tunnel_q9} and the aligned pathway's 5\% entrance. TimesFM-3 uses the zero-anchored adapter; the identity-anchored adapter of the submission gives 4 / 5 / 9 and $0.05$.}
\vspace{4pt}
\label{tab:native_reference}
\small
\setlength{\tabcolsep}{4pt}
\begin{tabular}{@{}lcccc@{}}
\toprule
& \multicolumn{1}{c}{Final-block readout} & \multicolumn{1}{c}{Readout} & \multicolumn{1}{c}{Aligned pathway} & \\
\cmidrule(lr){2-2}\cmidrule(lr){3-3}\cmidrule(lr){4-4}
Model & vs.\ full model (within 5\%) & 2\% / 5\% / 10\% & 2\% / 5\% / 10\% & Spearman \\
\midrule
Chronos-2 & $+17.7\%$ (2/14) & 0 / 2 / 4 & 8 / 13 / 14 & $0.81$ \\
TimesFM-3 & $+20.3\%$ (1/14) & 0 / 1 / 5 & 8 / 11 / 11 & $-0.02$ \\
TiRex & $+12.7\%$ (3/14) & 1 / 3 / 6 & 8 / 9 / 14 & $0.80$ \\
\midrule
All & 6/42 & 1 / 6 / 15 & 24 / 33 / 39 & \\
\bottomrule
\end{tabular}
\end{table}



\subsection{Additional geometry results}
\label{app:geometry}

\paragraph{Geometry estimators.}
We compute geometry from the same forecast-producing states used by the readouts, using the
training split. Chronos-2 contributes its four forecast-slot states per window, TimesFM-3 its last
context token, and TiRex the corresponding state from the first inference pass. We use debiased
linear CKA \citep{kornblith2019similarityneuralnetworkrepresentations,
cortes2024algorithmslearningkernelsbased,murphy2024correctingbiasedcenteredkernel} and effective
rank \citep{7098875}.
Training-split CKA provides substantially more samples than the test split; as a robustness check,
the dominant CKA boundary differs by at most one block between the two splits in 37 of 42
model-dataset pairs. Chronos-2's L0 forecast-slot states contain no window-specific information,
so CKA involving L0 is excluded.

\paragraph{Geometric landmarks.}
We summarize each profile by the effective-rank peak, the dominant CKA boundary, the change in
CKA and effective rank at the tunnel entrance, and the sharpest adjacent CKA change. The dominant
CKA boundary is the split of L1 to $\mathrm{L}L$ into two contiguous regions that maximizes
within-region relative to between-region CKA. A landmark is considered near the tunnel entrance
if it lies within $\max(1,\operatorname{round}(0.1L))$ blocks. These descriptors are post hoc
summaries rather than selection criteria.

\begin{table}[t]
  \centering
  \caption{\textbf{Forecast recoverability versus representation geometry.} Counts of
  model-dataset pairs. Entry-step statistics exclude the seven pairs whose tunnel begins at L0.
  Chance denotes the expected count under a uniformly placed landmark independent of the entrance.
  The 2\% columns repeat the analysis with the entrances of the stricter 2\% tolerance; entry-step
  statistics then exclude the one pair whose tunnel begins at L0.}
  \vspace{4pt}
  \label{tab:geometry_summary}
  \small
  \setlength{\tabcolsep}{5pt}
  \begin{tabular}{@{}lcccc@{}}
  \toprule
  & \multicolumn{2}{c}{5\% tunnel} & \multicolumn{2}{c}{2\% tunnel} \\
  \cmidrule(lr){2-3}\cmidrule(l){4-5}
  Observation & Pairs & Chance & Pairs & Chance \\
  \midrule
  Entry step is the largest adjacent CKA change & 3/35 & 2.6 & 4/41 & 3.0 \\
  Entry step is the largest adjacent effective-rank change & 2/35 & 2.6 & 2/41 & 3.0 \\
  Effective-rank peak is near the entrance & 5/35 & 7.8 & 10/41 & 9.0 \\
  Dominant CKA boundary is near the entrance & 8/35 & 7.8 & 12/41 & 9.0 \\
  \midrule
  Sharpest adjacent CKA change lies inside the tunnel & 34/42 & 24.4 & 30/42 & 16.1 \\
  Sharpest adjacent CKA change is the step into the final block & 25/42 & & 25/42 & \\
  Effective-rank peak lies at or after the entrance & 29/42 & 24.5 & 23/42 & 16.9 \\
  At least half of the effective-rank variation lies inside the tunnel & 28/42 & & 20/42 & \\
  CKA between the entrance and final block is below 0.5 & 23/42 & & 16/42 & \\
  \bottomrule
  \end{tabular}
\end{table}

The geometric landmarks do not consistently coincide with the forecasting tunnel entrance
(Table~\ref{tab:geometry_summary}). Instead, substantial geometric change often continues inside
the tunnel: the sharpest adjacent CKA change lies within it in 34 of 42 model-dataset pairs.
The final block retains a median 7\%, 19\%, and 9\% of peak effective rank for Chronos-2,
TimesFM-3, and TiRex, respectively. The same pattern holds under the stricter 2\% tolerance, whose
entrances lie at or after the 5\% entrances: the sharpest adjacent CKA change lies inside the tunnel
in 30 of 42 pairs (16.1 expected by chance), and the effective-rank peak lies at or after the
entrance in 23 of 42 (16.9 expected). Because a larger tolerance can only move the entrance
earlier, both counts can only grow for tolerances above 5\%. Thus, forecasting performance can
saturate well before representation geometry does.

Figure~\ref{fig:app_cka_chronos2} shows the full Chronos-2 CKA matrices. Several datasets display
blocks of mutually similar depths separated by sharper transitions, so CKA suggests several
candidate cut points rather than a single one, and none of them is tied to forecast quality.

\begin{figure}[p]
  \centering
  \includegraphics[width=\textwidth]{figures/appendix_cka_chronos2.pdf}
  \caption{\textbf{Chronos-2 debiased linear CKA across all depths and datasets.}
  CKA is computed on training representations; dashed lines mark the validation-selected
  sustained 5\% tunnel entrance.}
  \label{fig:app_cka_chronos2}
\end{figure}

\subsection{Additional compatibility results}
\label{app:compatibility}

Figures~\ref{fig:app_h3_depth_chronos2}, \ref{fig:app_h3_depth_timesfm3} and
\ref{fig:app_h3_depth_tirex} show the test MASE increase of every readout at every depth for all
14 datasets, in the format of Figure~\ref{fig:h3_main}. They also show a forecast-trained adapter
for Chronos-2 and TiRex: the same affine map as in Section~\ref{sec:adapter}, trained instead on the
$Q=9$ pinball loss against the true future through the frozen pathway, with the iterative protocol
of Appendix~\ref{app:alignment}. At the tunnel entrance its median test MASE increase is no lower
than the label-free adapter's ($+6.7\%$ against $+5.5\%$ for Chronos-2, $+8.7\%$ against $+6.3\%$
for TiRex), although it is slightly lower at a few early depths (Chronos-2 L1; TiRex L1, L3 and
L4).

\paragraph{Depth profiles.} Aligned Chronos-2 has a median test MASE increase between $+0.4\%$
(L11) and $+4.4\%$ (L3) at every depth from L3 onwards. For TimesFM-3, the zero-anchored aligned
pathway stays between $+6\%$ and $+12\%$ from L1 to L17 and reaches $+3.9\%$ at L18 and $+1.3\%$ at
L19; the identity-anchored adapter of the submitted version stayed between roughly $+30\%$ and
$+82\%$ up to L17 (Appendix~\ref{sec:app_timesfm3_ridge}).

\paragraph{Nonlinear adapter.} Table~\ref{tab:h3_depth_nonlinear} repeats Table~\ref{tab:h3_depth}
with the nonlinear adapter of Appendix~\ref{app:alignment}. For Chronos-2 and TiRex, the median
change in test MASE over all 168 depth-dataset pairs of each model is $-0.2$ percentage points, and
the share of the hard-cut gap closed moves from 85 to 87\% (Chronos-2) and from 90 to 91\% (TiRex).
The affine map therefore captures almost all of the compatibility that a lightweight adapter can
restore. For TimesFM-3 the nonlinear adapter is far worse than the affine one, because it must be
fitted iteratively and this fit does not reproduce the full model's forecasts even on the training
windows (Appendix~\ref{app:alignment}).

\begin{table}[h]
  \centering
  \caption{\textbf{Affine versus nonlinear adapter at 25, 50 and 75\% of network depth.} Median
  test MASE increase over the full model across all 14 datasets, with the number of datasets within
  5\% in parentheses. TimesFM-3 affine: the zero-anchored closed form
  (Appendix~\ref{sec:app_timesfm3_ridge}). $^\dagger$Iterative fit; on Electricity it does not
  reproduce the full model's forecasts even on the training windows (training $R^2<0$).}
  \vspace{4pt}
  \label{tab:h3_depth_nonlinear}
  \small
  \setlength{\tabcolsep}{4pt}
  \begin{tabular}{llrr}
  \toprule
  Model & Depth & Affine & Nonlinear \\
  \midrule
  Chronos-2 & L3 (25\%)  & $+4.4\%$ (7)  & $+3.6\%$ (8)  \\
            & L6 (50\%)  & $+2.3\%$ (9)  & $+2.2\%$ (9)  \\
            & L9 (75\%)  & $+2.2\%$ (10) & $+1.6\%$ (11) \\
  \midrule
  TimesFM-3 & L5 (25\%)  & $+10.6\%$ (4) & $7.1\times^\dagger$ (0)  \\
            & L10 (50\%) & $+11.3\%$ (2) & $14.5\times^\dagger$ (0) \\
            & L15 (75\%) & $+8.5\%$ (3) & $20.5\times^\dagger$ (0) \\
  \midrule
  TiRex     & L3 (25\%)  & $+13.4\%$ (1) & $+15.3\%$ (2) \\
            & L6 (50\%)  & $+10.6\%$ (3) & $+8.7\%$ (5)  \\
            & L9 (75\%)  & $+6.2\%$ (5)  & $+6.2\%$ (5)  \\
  \bottomrule
  \end{tabular}
\end{table}

\begin{figure}[p]
  \centering
  \includegraphics[width=\textwidth]{figures/h3_depth_chronos2.pdf}
  \caption{\textbf{Chronos-2: native-pathway compatibility across depth} on all 14 datasets. Test
  MASE increase over the full model for the hard cut, the aligned native pathway (with its 95\%
  series-level cluster-bootstrap band), the forecast-trained adapter and the probe head
  (reference, dashed). The shaded region is the forecasting tunnel.}
  \label{fig:app_h3_depth_chronos2}
\end{figure}

\begin{figure}[p]
  \centering
  \includegraphics[width=\textwidth]{figures/h3_depth_timesfm3_rev.pdf}
  \caption{\textbf{TimesFM-3: native-pathway compatibility across depth} on all 14 datasets, in
  the format of Figure~\ref{fig:app_h3_depth_chronos2}. The aligned pathway uses the zero-anchored
  adapter; the dotted line is the identity-anchored adapter of the submitted version
  (Appendix~\ref{sec:app_timesfm3_ridge}). A forecast-trained adapter is not shown: with TimesFM-3's
  linear head it spans exactly the function class of the probe head.}
  \label{fig:app_h3_depth_timesfm3}
\end{figure}

\begin{figure}[p]
  \centering
  \includegraphics[width=\textwidth]{figures/h3_depth_tirex.pdf}
  \caption{\textbf{TiRex: native-pathway compatibility across depth} on all 14 datasets, in the
  format of Figure~\ref{fig:app_h3_depth_chronos2}.}
  \label{fig:app_h3_depth_tirex}
\end{figure}

% Appendix B.3.1: regularization of the closed-form TimesFM-3 adapter (Phase 2g, review period).
% Numbers: results/three_model_phase2g_timesfm3_zero_ridge/ (Narval job 9588346) and
% experiments/make_rebuttal_numbers.py, section D. Figure: experiments/make_rebuttal_figures.py.

\subsubsection{Regularization of the TimesFM-3 adapter}
\label{sec:app_timesfm3_ridge}

The submitted closed-form TimesFM-3 adapter penalized $\lambda\|\Delta_\ell\|_F^2$, i.e.\ it shrank $A_\ell=I+\Delta_\ell$ toward the identity and therefore toward the hard cut. For TimesFM-3 the hard cut is 35 to 160 times worse than the full model in median test MASE from L0 to L17 (Table~\ref{tab:timesfm3_ridge}). Validation therefore selected almost no regularization, $\kappa\le10^{-8}$ in 231 of the 280 fits (14 datasets $\times$ 20 depths), and the adapter nearly interpolated its 1,394 training windows: from L1 to L17, its median validation distillation error was 13 to 184 times its training error, and its test $R^2$ against the full model's forecast was negative on 4 to 11 of the 14 datasets (Table~\ref{tab:timesfm3_ridge_fit}). At L0 the same fit retained far fewer effective parameters and generalized, which is why aligned TimesFM-3 at L0 outperformed L2 to L17 in the submission.

\paragraph{Zero anchor.} We refit the adapter with the penalty on the whole linear map, $\lambda\|I+\Delta_\ell\|_F^2$ (Equation~\ref{eq:adapter_ridge}), which shrinks $A_\ell$ toward zero, and changed nothing else: the same cached states, splits, teacher forecasts, metric $W^\top W$, $\lambda=\kappa\,g_{\max}m_{\max}$ with $\kappa\in\{10^{-12},\dots,10^{2}\}$, and selection on validation distillation error, with ties going to the larger $\lambda$. The limit $\lambda\to\infty$, $A_\ell=0$ with $b_\ell$ the mean teacher state, which produces a constant forecast, takes the place of the hard cut as the explicit candidate. Penalizing $\|W A_\ell\|_F^2$ instead, i.e.\ plain ridge regression of the full model's head outputs on the states, gives similar results (Table~\ref{tab:timesfm3_ridge}, last aligned column), so the conclusions do not depend on the form of the penalty. The identity-anchored fit is reproduced exactly by the same code: its test MASE matches the submission to within $3.4\times10^{-5}$ relative at every depth and dataset, with the same selected $\kappa$ in all 280 fits.

\paragraph{Results.} The zero anchor changes the aligned TimesFM-3 pathway substantially (Table~\ref{tab:timesfm3_ridge}, Figure~\ref{fig:app_h3_depth_timesfm3}). From L1 to L14 its median test MASE increase stays between $+9.9\%$ and $+11.9\%$, against $+38.7\%$ to $+81.9\%$ with the identity anchor, and it is better on all 14 datasets at L1 to L8 and L10 to L15 (paired 95\% intervals). L0 becomes its worst depth: the zero-anchored pathway at L0 is worse than at every depth from L1 to L18 on 12 of the 14 datasets. Validation alone prefers the zero anchor in all 280 fits. Its median validation error is 1.9 to 3.2 times its training error, comparable to the iteratively fitted Chronos-2 and TiRex adapters (median 1.8 and 2.3 at the selected checkpoint).

The corrected aligned pathway also has a lower test MASE than the supervised readout on 10 to 12 of the 14 datasets at every depth from L1 (paired 95\% intervals). Both are affine maps of the same state, but the readout is fitted to the observed future with the pinball loss, the adapter to the full model's forecasts; we do not attribute the difference to a single cause. The exceptions are BOOM and Wind Farms, where validation selects $\kappa\ge10$, i.e.\ a nearly constant forecast, at L1 to L18; there, none of our affine fits reproduces the full model's forecasts on validation windows, and the aligned pathway remains far behind the readout ($+181\%$ and $+38\%$ at L5, against $+47\%$ and $+10\%$). On LOOP Seattle the readout is 8\% to 20\% better than the full model itself, which a fit to the full model's forecasts cannot follow.

\paragraph{Order of analysis and physical truncation.} We introduced the zero anchor after the test results of the identity-anchored fit had been computed, prompted by its anomalous L0 performance; both fits are reported here, and the main text uses the zero-anchored one. All TimesFM-3 accuracy numbers with the zero anchor are computed offline, from cached representations. The adapter has the same shape as before, so the measured latency is unchanged. We repeated the physical check of Appendix~\ref{app:latency} for the zero-anchored adapters at L0, L1, L5, L10, L15, L18, and L19 on all 14 datasets, 98 configurations evaluated on CPUs. Each adapter is refitted from the cached states with the same code, selects the same $\kappa$, and reproduces the offline numbers exactly; it is then spliced before the head of a copy of the public TimesFM-3 model with the blocks after L$\ell$ removed. In every configuration, the physically truncated model's mean test MASE agrees with the offline pathway to within $4.5\times10^{-5}$ relative, satisfying the rule of Appendix~\ref{app:latency}, and the unmodified model agrees with the cached full-model forecasts to within $4.2\times10^{-5}$. As for the original adapters, individual forecast values agree less tightly: the median difference is about $10^{-6}$ of the mean forecast magnitude, and the largest is 14\% on BOOM at L0 and at most 2.3\% elsewhere.

\begin{table}[t]
\centering
\caption{\textbf{TimesFM-3 aligned pathway with identity- and zero-anchored regularization.} Median test MASE increase over the full model across 14 datasets; $\times$ denotes ratios to the full-model MASE. Zero, $\|WA\|$: the zero-anchored fit with the penalty $\lambda\|W A_\ell\|_F^2$. $\le$5\%: datasets within 5\% of the full model (zero anchor). Zero $<$ identity and zero vs.\ readout: datasets where the paired 95\% interval of the ratio excludes one, better / worse. Speedup: end to end at batch size 1, from the physically truncated identity-anchored models; the zero-anchored adapter has the same shape.}
\vspace{4pt}
\label{tab:timesfm3_ridge}
\small
\setlength{\tabcolsep}{3.5pt}
\begin{tabular}{@{}lrrrrrcccr@{}}
\toprule
& & \multicolumn{3}{c}{Aligned pathway} & & & & & \\
\cmidrule(lr){3-5}
Cut & Hard cut & Identity & Zero & Zero, $\|WA\|$ & Readout & $\le$5\% & Zero $<$ id. & Zero vs.\ readout & Speedup \\
\midrule
L0 & $35\times$ & $+31.3\%$ & $+24.3\%$ & $+23.9\%$ & $+31.9\%$ & 0 & 8 & 8 / 2 & $1.38\times$ \\
L1 & $38\times$ & $+38.7\%$ & $+10.6\%$ & $+9.7\%$ & $+15.8\%$ & 2 & 14 & 10 / 2 & $1.37\times$ \\
L5 & $53\times$ & $+69.8\%$ & $+10.6\%$ & $+8.5\%$ & $+16.8\%$ & 4 & 14 & 10 / 4 & $1.24\times$ \\
L10 & $121\times$ & $+67.3\%$ & $+11.3\%$ & $+10.1\%$ & $+17.0\%$ & 2 & 14 & 10 / 3 & $1.19\times$ \\
L15 & $160\times$ & $+47.9\%$ & $+8.5\%$ & $+7.1\%$ & $+14.8\%$ & 3 & 14 & 10 / 3 & $1.08\times$ \\
L17 & $91\times$ & $+29.7\%$ & $+6.1\%$ & $+6.0\%$ & $+14.7\%$ & 5 & 13 & 11 / 3 & $1.07\times$ \\
L18 & $62\times$ & $+14.3\%$ & $+3.9\%$ & $+3.7\%$ & $+13.2\%$ & 9 & 13 & 11 / 3 & $1.03\times$ \\
L19 & $7\times$ & $+5.4\%$ & $+1.3\%$ & $+1.0\%$ & $+12.7\%$ & 11 & 13 & 12 / 1 & $0.99\times$ \\
\bottomrule
\end{tabular}
\end{table}

\begin{table}[t]
\centering
\caption{\textbf{What validation selected for the two TimesFM-3 adapters.} Medians over the 14 datasets at each depth, summarized over L1 to L17; L0 in parentheses. Effective parameters: the trace of the ridge smoother for each of the head's output directions, averaged with weights proportional to the full model's output variance in that direction (at most 1,280). Validation / training: distillation error of the selected fit on the validation windows relative to the training windows. Test $R^2$: against the full model's forecast in the head's output metric.}
\vspace{4pt}
\label{tab:timesfm3_ridge_fit}
\small
\setlength{\tabcolsep}{5pt}
\begin{tabular}{@{}lcc@{}}
\toprule
& Identity anchor (submission) & Zero anchor \\
\midrule
Selected $\kappa$ & $\le10^{-8}$ in 231 of 280 fits & $10^{-6}$ to $10^{-4}$ in 221 of 280 fits \\
Effective parameters & 1,017 to 1,280 (253) & 262 to 681 (76) \\
Validation / training error & 13 to 184 (1.9) & 1.9 to 3.2 (1.3) \\
Test $R^2$ & $-0.60$ to $0.53$ (0.12) & $0.64$ to $0.88$ (0.33) \\
Datasets with test $R^2<0$ & 4 to 11 (6) & 2 to 3 (3) \\
Lower validation error & 0 of 280 fits & 280 of 280 fits \\
\bottomrule
\end{tabular}
\end{table}



\subsection{Additional truncation results}
\label{app:truncation_results}

\paragraph{Complete frontier.} Table~\ref{tab:h4_frontier_full} gives the complete numerical
version of the performance-compute frontier in Figure~\ref{fig:h4_frontier}: every aligned
truncation depth of all three models, with the full model and Chronos-2-small as references.
Forecasting performance is computed as described in Appendix~\ref{app:latency}.

\paragraph{Throughput at batch size 256.} For Chronos-2 and TiRex, speedups at batch size 256
closely track those at batch size 1; for example, Chronos-2 after L4 runs $2.43\times$ and
$2.51\times$ faster, and TiRex after L6 $1.99\times$ and $1.96\times$. TimesFM-3 differs: its
speedup at batch size 1 never exceeds $1.43\times$, whereas at batch size 256 it grows steadily
with the number of removed blocks, to $2.70\times$ after L5 and $5.94\times$ with no block retained.
This is consistent with a fixed per-call cost outside the backbone that dominates at batch size 1
(Appendix~\ref{app:latency}). Its forecasting degradation does not depend on the batch size, so
with the zero-anchored adapter truncation after L1 runs $4.65\times$ faster at batch size 256 for a
median $+10.6\%$, and after L5 $2.70\times$ faster for the same increase. Chronos-2-small also
benefits more from batching ($1.80\times$ at batch size 1, $3.28\times$ at 256).

\begin{table}[p]
\centering
\caption{\textbf{Performance-compute frontier of aligned truncation, all depths.} ``L$k$'' keeps
blocks $1,\dots,k$ and removes the rest; L0 keeps no block. MASE increase is the median over the
14 datasets of test MASE relative to the full model, with its interquartile range. Speedups are
end to end, from a context in host memory to forecasts returned to host memory, relative to the
full model at the same batch size: batch size 1 (B=1) measures single-forecast latency and is the
headline setting; batch size 256 (B=256) measures throughput. Chronos-2-small is a separately
pretrained model, not a truncation of Chronos-2. TimesFM-3 uses the zero-anchored adapter, with
accuracy computed offline (Appendix~\ref{sec:app_timesfm3_ridge}).}
\vspace{4pt}
\label{tab:h4_frontier_full}
\small
\setlength{\tabcolsep}{5pt}
\begin{tabular}{@{}llrcrr@{}}
\toprule
Model & Cut & \multicolumn{2}{c}{Test MASE increase} & \multicolumn{2}{c}{Speedup} \\
\cmidrule(lr){3-4}\cmidrule(l){5-6}
 & & Median & IQR & B=1 & B=256 \\
\midrule
Chronos-2 & L12 (full model) & $0.0\%$ &  & $1.00\times$ & $1.00\times$ \\
 & L11 & $+0.4\%$ & [$0.0\%$, $+1.2\%$] & $1.08\times$ & $1.08\times$ \\
 & L10 & $+0.9\%$ & [$+0.2\%$, $+3.3\%$] & $1.18\times$ & $1.17\times$ \\
 & L9 & $+2.2\%$ & [$+0.9\%$, $+6.2\%$] & $1.28\times$ & $1.29\times$ \\
 & L8 & $+3.1\%$ & [$+2.0\%$, $+8.2\%$] & $1.42\times$ & $1.43\times$ \\
 & L7 & $+3.1\%$ & [$+0.8\%$, $+9.6\%$] & $1.57\times$ & $1.60\times$ \\
 & L6 & $+2.3\%$ & [$+0.4\%$, $+11.5\%$] & $1.78\times$ & $1.82\times$ \\
 & L5 & $+2.1\%$ & [$+0.5\%$, $+12.1\%$] & $2.06\times$ & $2.11\times$ \\
 & L4 & $+3.1\%$ & [$+0.5\%$, $+11.4\%$] & $2.43\times$ & $2.51\times$ \\
 & L3 & $+4.4\%$ & [$+0.7\%$, $+11.3\%$] & $2.91\times$ & $3.10\times$ \\
 & L2 & $+8.4\%$ & [$+2.1\%$, $+30.0\%$] & $3.72\times$ & $4.07\times$ \\
 & L1 & $+13.1\%$ & [$+6.8\%$, $+51.9\%$] & $5.19\times$ & $5.88\times$ \\
 & L0 & $+69.4\%$ & [$+19.9\%$, $+183.8\%$] & $8.41\times$ & $9.93\times$ \\
\addlinespace
Chronos-2-small & separate model & $+0.8\%$ & [$0.0\%$, $+5.9\%$] & $1.80\times$ & $3.28\times$ \\
\midrule
TimesFM-3 & L20 (full model) & $0.0\%$ &  & $1.00\times$ & $1.00\times$ \\
 & L19 & $+1.3\%$ & [$+0.6\%$, $+3.9\%$] & $0.99\times$ & $1.04\times$ \\
 & L18 & $+3.9\%$ & [$+1.6\%$, $+11.4\%$] & $1.03\times$ & $1.09\times$ \\
 & L17 & $+6.1\%$ & [$+2.8\%$, $+19.3\%$] & $1.07\times$ & $1.15\times$ \\
 & L16 & $+7.1\%$ & [$+4.1\%$, $+23.6\%$] & $1.05\times$ & $1.20\times$ \\
 & L15 & $+8.5\%$ & [$+6.0\%$, $+24.4\%$] & $1.08\times$ & $1.26\times$ \\
 & L14 & $+9.9\%$ & [$+5.4\%$, $+25.7\%$] & $1.11\times$ & $1.34\times$ \\
 & L13 & $+10.4\%$ & [$+5.5\%$, $+27.7\%$] & $1.11\times$ & $1.42\times$ \\
 & L12 & $+11.1\%$ & [$+6.1\%$, $+31.5\%$] & $1.15\times$ & $1.51\times$ \\
 & L11 & $+11.2\%$ & [$+6.4\%$, $+31.2\%$] & $1.12\times$ & $1.60\times$ \\
 & L10 & $+11.3\%$ & [$+6.6\%$, $+31.8\%$] & $1.19\times$ & $1.71\times$ \\
 & L9 & $+11.9\%$ & [$+6.4\%$, $+32.5\%$] & $1.18\times$ & $1.85\times$ \\
 & L8 & $+11.3\%$ & [$+6.5\%$, $+36.7\%$] & $1.25\times$ & $2.03\times$ \\
 & L7 & $+10.5\%$ & [$+6.4\%$, $+35.4\%$] & $1.22\times$ & $2.21\times$ \\
 & L6 & $+10.4\%$ & [$+7.0\%$, $+31.7\%$] & $1.24\times$ & $2.44\times$ \\
 & L5 & $+10.6\%$ & [$+4.5\%$, $+31.3\%$] & $1.24\times$ & $2.70\times$ \\
 & L4 & $+11.3\%$ & [$+4.9\%$, $+31.3\%$] & $1.29\times$ & $3.11\times$ \\
 & L3 & $+11.0\%$ & [$+4.8\%$, $+30.7\%$] & $1.34\times$ & $3.49\times$ \\
 & L2 & $+11.1\%$ & [$+5.1\%$, $+31.5\%$] & $1.43\times$ & $4.25\times$ \\
 & L1 & $+10.6\%$ & [$+7.2\%$, $+32.4\%$] & $1.37\times$ & $4.65\times$ \\
 & L0 & $+24.3\%$ & [$+15.7\%$, $+48.8\%$] & $1.38\times$ & $5.94\times$ \\
\midrule
TiRex & L12 (full model) & $0.0\%$ &  & $1.00\times$ & $1.00\times$ \\
 & L11 & $+1.8\%$ & [$+0.8\%$, $+5.3\%$] & $1.09\times$ & $1.09\times$ \\
 & L10 & $+4.3\%$ & [$+2.0\%$, $+6.2\%$] & $1.21\times$ & $1.19\times$ \\
 & L9 & $+6.2\%$ & [$+3.4\%$, $+10.4\%$] & $1.33\times$ & $1.32\times$ \\
 & L8 & $+8.7\%$ & [$+4.8\%$, $+12.8\%$] & $1.50\times$ & $1.48\times$ \\
 & L7 & $+8.5\%$ & [$+6.7\%$, $+13.2\%$] & $1.71\times$ & $1.69\times$ \\
 & L6 & $+10.6\%$ & [$+5.7\%$, $+14.7\%$] & $1.99\times$ & $1.96\times$ \\
 & L5 & $+11.3\%$ & [$+6.6\%$, $+20.1\%$] & $2.39\times$ & $2.33\times$ \\
 & L4 & $+12.8\%$ & [$+7.4\%$, $+28.2\%$] & $2.99\times$ & $2.88\times$ \\
 & L3 & $+13.4\%$ & [$+8.2\%$, $+33.8\%$] & $3.93\times$ & $3.76\times$ \\
 & L2 & $+10.5\%$ & [$+7.9\%$, $+36.1\%$] & $5.95\times$ & $5.43\times$ \\
 & L1 & $+14.9\%$ & [$+8.1\%$, $+44.5\%$] & $11.7\times$ & $9.83\times$ \\
 & L0 & $+46.2\%$ & [$+11.9\%$, $+109.2\%$] & $340\times$ & $50.2\times$ \\
\bottomrule
\end{tabular}
\end{table}

\paragraph{Truncation inside and before the tunnel.} Table~\ref{tab:h4_tunnel_strata} gives the
depth-matched comparison of Section~\ref{sec:results_acceleration}. At each cut depth, the datasets
of a model are split according to whether that depth lies inside their validation-selected 5\%
tunnel ($\ell\geq\ell_{\mathrm{tunnel}}$) or before it, and the median aligned test MASE increase
is reported for each group. Because speedup depends only on the cut depth, both groups are compared
at the same speedup. Depths are listed where both groups contain at least two datasets. Truncation
is cheaper inside the tunnel at all 9 such depths for Chronos-2, all 12 for TiRex, and all 15 for
TimesFM-3 with the zero-anchored adapter (14 of 15 with the identity-anchored adapter of the
submitted version, the exception being L13, where only three datasets lie before the tunnel).
The comparison is observational: datasets whose tunnel begins early may also be easier to forecast
from intermediate states for other reasons.

\begin{table}[p]
\centering
\caption{\textbf{Aligned truncation inside versus before the forecasting tunnel.} Median test MASE
increase over the full model, at the same cut depth, for datasets where that depth lies inside
their validation-selected 5\% tunnel and for datasets where it lies before the tunnel entrance.
$n$ gives the number of datasets in each group. TimesFM-3 uses the zero-anchored adapter
(Appendix~\ref{sec:app_timesfm3_ridge}).}
\vspace{4pt}
\label{tab:h4_tunnel_strata}
\small
\setlength{\tabcolsep}{5pt}
\begin{tabular}{@{}llccc@{}}
\toprule
Model & Cut & $n$ (inside / before) & Inside tunnel & Before tunnel \\
\midrule
Chronos-2 & L1 & 2 / 12 & $+4.8\%$ & $+21.5\%$ \\
 & L2 & 4 / 10 & $+1.8\%$ & $+15.1\%$ \\
 & L3 & 7 / 7 & $+2.6\%$ & $+9.3\%$ \\
 & L4 & 8 / 6 & $+2.2\%$ & $+14.2\%$ \\
 & L5 & 8 / 6 & $+1.9\%$ & $+13.8\%$ \\
 & L6 & 8 / 6 & $+1.9\%$ & $+13.1\%$ \\
 & L7 & 8 / 6 & $+1.8\%$ & $+11.8\%$ \\
 & L8 & 9 / 5 & $+2.4\%$ & $+8.3\%$ \\
 & L9 & 11 / 3 & $+1.8\%$ & $+4.5\%$ \\
\midrule
TimesFM-3 & L0 & 4 / 10 & $+16.5\%$ & $+29.5\%$ \\
 & L1 & 8 / 6 & $+10.6\%$ & $+15.0\%$ \\
 & L2 & 8 / 6 & $+11.1\%$ & $+12.1\%$ \\
 & L3 & 8 / 6 & $+11.0\%$ & $+11.5\%$ \\
 & L4 & 8 / 6 & $+11.3\%$ & $+12.6\%$ \\
 & L5 & 9 / 5 & $+10.3\%$ & $+19.5\%$ \\
 & L6 & 9 / 5 & $+10.3\%$ & $+18.7\%$ \\
 & L7 & 9 / 5 & $+10.3\%$ & $+19.1\%$ \\
 & L8 & 9 / 5 & $+10.5\%$ & $+20.9\%$ \\
 & L9 & 9 / 5 & $+10.3\%$ & $+21.1\%$ \\
 & L10 & 9 / 5 & $+10.6\%$ & $+22.7\%$ \\
 & L11 & 9 / 5 & $+10.3\%$ & $+22.6\%$ \\
 & L12 & 11 / 3 & $+10.9\%$ & $+13.8\%$ \\
 & L13 & 11 / 3 & $+10.4\%$ & $+13.0\%$ \\
 & L14 & 12 / 2 & $+8.9\%$ & $+20.6\%$ \\
\midrule
TiRex & L0 & 2 / 12 & $+8.0\%$ & $+53.3\%$ \\
 & L1 & 4 / 10 & $+6.8\%$ & $+25.5\%$ \\
 & L2 & 4 / 10 & $+6.5\%$ & $+24.0\%$ \\
 & L3 & 4 / 10 & $+6.7\%$ & $+23.3\%$ \\
 & L4 & 4 / 10 & $+7.7\%$ & $+17.3\%$ \\
 & L5 & 4 / 10 & $+5.8\%$ & $+13.7\%$ \\
 & L6 & 4 / 10 & $+4.4\%$ & $+12.3\%$ \\
 & L7 & 4 / 10 & $+4.6\%$ & $+11.7\%$ \\
 & L8 & 6 / 8 & $+4.9\%$ & $+10.0\%$ \\
 & L9 & 7 / 7 & $+2.9\%$ & $+8.3\%$ \\
 & L10 & 9 / 5 & $+2.2\%$ & $+9.8\%$ \\
 & L11 & 12 / 2 & $+1.5\%$ & $+8.2\%$ \\
\bottomrule
\end{tabular}
\end{table}

% Appendix B.4.1: cutting at the tunnel entrance (pre-specified test, reported in the review period).
% Numbers: experiments/make_rebuttal_numbers.py, section B; physical evaluation:
% results/three_model_phase2/combined/h4_outcomes.csv.

\subsubsection{Cutting at the tunnel entrance}
\label{sec:app_entrance_rule}

Before computing the truncation results, we specified one test of the tunnel as a cut rule: truncate each model at its validation-selected 5\% entrance, evaluate no other depth, and count the datasets on which the aligned model stays within 5\% of the full model in test MASE. This test was evaluated on physically truncated models (Appendix~\ref{app:latency}), whose test MASE agrees with the offline pathway to within $10^{-4}$ relative. Table~\ref{tab:entrance_rule} compares it with a single fixed depth for all datasets at a similar speed.

The entrance is not a reliable cut rule on its own: it keeps the aligned model within 5\% of the full model on 7 of 14 datasets for Chronos-2 and 5 of 12 for TiRex, whose entrance is the final block on M4 and Monash Traffic. For TimesFM-3 it does so on none of the 14 datasets with the identity-anchored adapter of the submission and on 2 with the zero-anchored adapter of Appendix~\ref{sec:app_timesfm3_ridge}, whose numbers are offline and post hoc. Compared with a fixed depth, the entrance gives a similar or slightly higher median but roughly halves the worst case: $+15.8\%$ against $+32.9\%$ at L4 and $+33.3\%$ at L3 for Chronos-2, and $+13.6\%$ against $+27.3\%$ at L9 for TiRex. Because the entrance adapts the cut to the dataset, it protects against datasets that tolerate truncation poorly, such as M4 and Electricity, without locating the best cut. The aligned pathway's own compatibility depth is rank-correlated with the entrance for Chronos-2 and TiRex but not for TimesFM-3 (Appendix~\ref{sec:app_native_reference}).

\begin{table}[t]
\centering
\caption{\textbf{Cutting at the tunnel entrance versus at a fixed depth.} Aligned pathway, test MASE increase over the full model. Entrance: each dataset cut at its validation-selected 5\% tunnel entrance (TiRex: the 12 datasets whose entrance precedes the final block). Within 5\%: datasets within 5\% of the full model. Speedup: median end-to-end speedup at batch size 1. TimesFM-3 with the zero anchor: offline, analysed after the test results of the identity-anchored fit.}
\vspace{4pt}
\label{tab:entrance_rule}
\small
\setlength{\tabcolsep}{5pt}
\begin{tabular}{@{}llcccc@{}}
\toprule
Model & Cut & Median & Within 5\% & Worst & Speedup \\
\midrule
Chronos-2 & Entrance & $+5.5\%$ & 7/14 & $+15.8\%$ & $2.67\times$ \\
 & Fixed L3 & $+4.4\%$ & 7/14 & $+33.3\%$ & $2.91\times$ \\
 & Fixed L4 & $+3.1\%$ & 8/14 & $+32.9\%$ & $2.43\times$ \\
\midrule
TiRex & Entrance & $+6.3\%$ & 5/12 & $+13.6\%$ & $1.42\times$ \\
 & Fixed L9 & $+6.2\%$ & 5/14 & $+27.3\%$ & $1.33\times$ \\
\midrule
TimesFM-3 & Entrance, identity anchor & $+35.5\%$ & 0/14 & $+293.0\%$ & $1.37\times$ \\
 & Entrance, zero anchor & $+14.2\%$ & 2/14 & $+83.1\%$ & $1.37\times$ \\
\bottomrule
\end{tabular}
\end{table}


% Appendix B.4.2: leave-one-dataset-out adapter transfer for Chronos-2 (Phase 2c, review period).
% Numbers: results/three_model_phase2_lodo/combined/ (Narval job 9229954). Figure generated by
% experiments/make_phase2_lodo_tables.py.

\subsubsection{Adapter transfer to unseen datasets}
\label{sec:app_lodo}

The adapters of the main text are fitted per dataset on unlabeled windows. To test whether this calibration can be avoided, we fit one adapter per depth on all datasets but the target and evaluate it once on the target, for Chronos-2 at L4, L6, and L9 ($2.43\times$, $1.78\times$, and $1.28\times$ faster at batch size 1). Electricity and Electricity (15-minute) share most of their meters, so when either is the target the other is also excluded, leaving 12 source datasets instead of 13. The objective, optimizer, weight-decay grid, and checkpoint selection are those of Appendix~\ref{app:alignment}. The objective and the validation error are averaged over source datasets with equal weight, each source contributing its 1,394 training windows. The target's windows are never used for fitting or selection. The pooled adapter has the same shape as a dataset-specific one, so the speedups carry over.

Table~\ref{tab:lodo} and Figure~\ref{fig:app_lodo} summarize the results. The transferred adapter is better than the hard cut on 13 of 14 datasets at each depth (paired 95\% intervals) and recovers a median 75\%, 70\%, and 80\% of the hard-cut gap at L4, L6, and L9, against 91\%, 86\%, and 87\% for the dataset-specific adapter. It is a median 7.1, 5.1, and 1.7 percentage points worse than the dataset-specific adapter, so transfer is close to calibration at L9 but not at L4. The exception is LOOP Seattle, the only 5-minute dataset, where the transferred adapter is worse than the hard cut at L4 and L6 ($+27.2\%$ and $+26.2\%$, against $+5.6\%$ and $+10.4\%$). The largest absolute degradation is on M4 at L4 ($+128\%$, against $+738\%$ for the hard cut and $+32.9\%$ for the dataset-specific adapter).

On the four datasets documented as held out from Chronos-2's pretraining (LOOP Seattle, SZ Taxi, M5, and Monash Traffic), the target appears neither in the backbone's documented pretraining data nor in the adapter's training data. There the transferred adapter improves on the hard cut on SZ Taxi, M5, and Monash Traffic at every depth, and fails on LOOP Seattle. Including the Electricity sibling among the sources changes the transferred adapter's MASE on the two Electricity targets by at most 2.5 percentage points. Of the 42 transferred fits, 19 selected a checkpoint in the last 10\% of the 1,500 epochs and were therefore still improving slowly; the weight decay never reached the grid maximum.

\paragraph{Cost of calibration.} A dataset-specific adapter needs the full model's own forecasts on unlabeled target windows: one forward pass over 1,394 windows, about 1\,s for Chronos-2 at batch size 256 on an A100, followed by a fit of about one minute on the same GPU, including the weight-decay sweep.

\begin{table}[t]
\centering
\caption{\textbf{Leave-one-dataset-out adapter transfer for Chronos-2.} Test MASE increase over the full model, median [IQR] over datasets. Transferred: one adapter fitted on the other datasets; dataset-specific: the adapter of the main text. Better than hard: datasets where the transferred adapter's paired 95\% interval lies below the hard cut. Within 5\% / 10\%: datasets within that margin of the full model (transferred adapter). Held out: the four datasets documented as held out from Chronos-2's pretraining.}
\vspace{4pt}
\label{tab:lodo}
\small
\setlength{\tabcolsep}{3.5pt}
\resizebox{\textwidth}{!}{%
\begin{tabular}{@{}llrcccccc@{}}
\toprule
Datasets & Cut & Speedup & Hard cut & Transferred & Dataset-specific & Better than hard & Within 5\% / 10\% & Gap (pp) \\
\midrule
All (14) & L4 & $2.43\times$ & $+37.0\%$ [$+13.8$, $+112.8$] & $+12.7\%$ [$+4.9$, $+24.5$] & $+3.1\%$ [$+0.5$, $+11.4$] & 13/14 & 4 / 5 & $+7.1$ \\
 & L6 & $1.78\times$ & $+20.1\%$ [$+11.3$, $+48.0$] & $+9.0\%$ [$+4.8$, $+18.5$] & $+2.3\%$ [$+0.4$, $+11.5$] & 13/14 & 4 / 8 & $+5.1$ \\
 & L9 & $1.28\times$ & $+17.1\%$ [$+11.3$, $+35.8$] & $+4.0\%$ [$+3.2$, $+7.6$] & $+2.2\%$ [$+0.9$, $+6.2$] & 13/14 & 9 / 11 & $+1.7$ \\
\midrule
Held out (4) & L4 & $2.43\times$ & $+8.4\%$ [$+5.5$, $+41.9$] & $+15.5\%$ [$+2.8$, $+29.0$] & $+0.8\%$ [$-1.0$, $+7.9$] & 3/4 & 2 / 2 & $+4.7$ \\
 & L6 & $1.78\times$ & $+10.6\%$ [$+8.6$, $+20.6$] & $+15.3\%$ [$+3.2$, $+27.8$] & $+1.1\%$ [$-0.9$, $+8.3$] & 3/4 & 2 / 2 & $+4.0$ \\
 & L9 & $1.28\times$ & $+6.3\%$ [$+0.9$, $+34.8$] & $+1.8\%$ [$-0.3$, $+6.5$] & $+1.1\%$ [$-0.8$, $+4.9$] & 3/4 & 3 / 3 & $+1.8$ \\
\bottomrule
\end{tabular}%
}
\end{table}

\begin{figure}[t]
\centering
\includegraphics[width=0.9\textwidth]{figures/lodo_per_dataset.pdf}
\caption{\textbf{Leave-one-dataset-out transfer for Chronos-2, per dataset.} Test MASE change over the full model after truncation at L4, L6, and L9 for the hard cut, the transferred adapter, and the dataset-specific adapter, with paired 95\% cluster-bootstrap intervals. Datasets are grouped by documented pretraining status for Chronos-2: held out, not listed, listed. The axis is linear to $\pm10\%$ and logarithmic beyond; triangles mark clipped values.}
\label{fig:app_lodo}
\end{figure}


% Appendix B.4.3: transfer as a function of the number of source datasets (Phase 2f, review period).
% Numbers: results/three_model_phase2f_source_scaling/combined/ (Narval job 9253740). Figure
% generated by experiments/make_phase2f_tables.py.

\subsubsection{Number of source datasets}
\label{sec:app_source_diversity}

We repeat the transfer experiment of Appendix~\ref{sec:app_lodo} at L6 with $N\in\{1,4,8,13\}$ source datasets for four target datasets. The targets were fixed from metadata before any result was computed: no Electricity dataset, a full test split of 262 series, four distinct sampling frequencies and domains, and all three documented pretraining groups. This gives LOOP Seattle (5-minute, held out), London Smart Meters (30-minute, listed), BOOM (hourly, not listed), and M5 (daily, held out). For each target, three random orders of the 13 other datasets, fixed by seed, define nested source sets: the $N=4$ set contains the $N=1$ source, and so on, while $N=13$ is the full pool, shared by the three orders. The two Electricity datasets never both appear in a set with $N\le8$. The fitting procedure is that of Appendix~\ref{sec:app_lodo}, so 40 adapters are fitted in total.

Table~\ref{tab:source_diversity} and Figure~\ref{fig:app_source_diversity} report the results. A single source dataset is unreliable: across the 12 target-source pairs, the MASE increase ranges from $+5.0\%$ to $+50.7\%$, it is worse than the hard cut in 3 pairs (paired 95\% intervals), and on average it is no better than the hard cut ($+19.7\%$ against $+19.8\%$). Four sources reduce the mean to $+12.7\%$ and make the result much less dependent on which sources are chosen; 8 and 13 sources give $+13.2\%$ and $+12.8\%$. The dataset-specific adapter remains better on all four targets at $N=13$ (paired 95\% intervals), with a mean of $+3.2\%$. LOOP Seattle does not benefit from more sources: its transferred adapter is worse than the hard cut in 8 of its 10 configurations. Each added source also adds 1,394 training windows, so the number of sources and the amount of training data increase together; we did not run a control that holds the total number of windows fixed.

\begin{table}[t]
\centering
\caption{\textbf{Transfer at L6 as a function of the number of source datasets $N$.} Test MASE increase over the full model in percent, for Chronos-2 truncated after L6 ($1.78\times$). For $N=1$, $4$, and $8$, the mean over three source orders with the range in brackets; $N=13$ is a single source set. Hard cut and dataset-specific adapter for reference.}
\vspace{4pt}
\label{tab:source_diversity}
\small
\setlength{\tabcolsep}{3pt}
\begin{tabular}{@{}lrrrrrr@{}}
\toprule
Target & Hard cut & $N=1$ & $N=4$ & $N=8$ & $N=13$ & Specific \\
\midrule
LOOP Seattle & $+10.4$ & $+25.8$ {\scriptsize[$10.6$, $44.8$]} & $+18.5$ {\scriptsize[$9.6$, $23.3$]} & $+26.8$ {\scriptsize[$20.7$, $31.8$]} & $+26.2$ & $-3.6$ \\
London Smart Meters & $+16.6$ & $+13.3$ {\scriptsize[$12.5$, $13.8$]} & $+11.8$ {\scriptsize[$10.9$, $12.4$]} & $+10.7$ {\scriptsize[$10.0$, $12.0$]} & $+9.6$ & $+4.6$ \\
BOOM & $+41.5$ & $+33.6$ {\scriptsize[$16.8$, $50.7$]} & $+16.4$ {\scriptsize[$12.0$, $21.9$]} & $+11.3$ {\scriptsize[$10.3$, $13.2$]} & $+11.2$ & $+9.3$ \\
M5 & $+10.8$ & $+6.0$ {\scriptsize[$5.0$, $6.8$]} & $+4.0$ {\scriptsize[$3.6$, $4.4$]} & $+3.8$ {\scriptsize[$3.2$, $4.3$]} & $+4.3$ & $+2.3$ \\
\midrule
Mean & $+19.8$ & $+19.7$ & $+12.7$ & $+13.2$ & $+12.8$ & $+3.2$ \\
\bottomrule
\end{tabular}
\end{table}

\begin{figure}[t]
\centering
\includegraphics[width=0.6\textwidth]{figures/source_scaling_main.pdf}
\caption{\textbf{Transfer at L6 as a function of the number of source datasets.} Mean test MASE change over the full model across three source orders, with the range as a band; crosses: hard cut; squares: dataset-specific adapter. The axis is linear to $\pm10\%$ and logarithmic beyond.}
\label{fig:app_source_diversity}
\end{figure}


% Appendix B.4.4: non-contiguous block selection for Chronos-2 (Phase 2e, review period).
% Numbers: results/three_model_phase2e_noncontig/tables/ (Narval job 9229964). Figure generated
% by experiments/make_phase2e_tables.py.

\subsubsection{Non-contiguous block selection}
\label{sec:app_noncontig}

Prefix truncation keeps blocks $1,\dots,\ell$. We compare it with three non-contiguous or shifted selections of the same number $k\in\{4,6,9\}$ of the 12 Chronos-2 blocks, fixed before any result was computed (Table~\ref{tab:noncontig}): uniformly spaced blocks, early plus late blocks, and a contiguous middle window as a control. Retained blocks keep their weights and order; skipped blocks are removed from the encoder, whose blocks each contain their own positional encoding, so no other change is needed. Every configuration receives the same label-free affine adapter as in the main text (Appendix~\ref{app:alignment}), fitted on the states of the reduced encoder, and every physically reduced model reproduces its offline aligned forecast in mean test MASE to within $10^{-4}$. Speedups were measured with the protocol of Appendix~\ref{app:latency} in a separate job of three repetitions; at a matched $k$ they agree within 6\%. For the prefix with $k=4$, this job measures $2.36\times$, against $2.43\times$ in Table~\ref{tab:h4_operating_points}.

At strong compression the prefix is best: with 4 blocks it is better than all three alternatives on 10, 11, and 12 of the 14 datasets, respectively, and worse on none (paired 95\% intervals), and validation selects it on 13 of 14 datasets. With 6 blocks, the prefix and uniform spacing are comparable. With 9 blocks, keeping the last blocks and skipping interior ones is better: uniform spacing improves on the prefix on 7 datasets and is worse on 1, and choosing the selection per dataset on validation lowers the median increase by 0.9 percentage points. These selections need almost no adapter: without one they cost $+1.1\%$ (uniform) and $+2.1\%$ (early plus late), against $+17.1\%$ for the prefix.

Removing a single block without an adapter explains the pattern (Figure~\ref{fig:app_single_skip}). Removing the final block costs a median $+18.5\%$ and degrades all 14 datasets; removing block 1 costs $+5.2\%$, block 2 $+2.2\%$, and block 11 $+1.5\%$, while each of blocks 3 to 10 costs at most $+0.4\%$. Block importance in Chronos-2 is therefore non-monotonic in depth: the first two blocks and the last two matter most. The final block produces the states the pretrained head expects, which is why the prefix needs an adapter, and why at mild compression keeping it is better than replacing it with one. These are four predefined selections on one model, not a search over selections, and the adapter corrects only the last retained state.

\begin{table}[t]
\centering
\caption{\textbf{Prefix versus non-contiguous block selection for Chronos-2.} Median test MASE increase over the full model across 14 datasets, aligned pathway (hard cut in parentheses). Better / worse: datasets where the selection is better / worse than the prefix with the same $k$ (aligned pathway, paired 95\% interval). Speedup: end to end at batch size 1, aligned pathway.}
\vspace{4pt}
\label{tab:noncontig}
\small
\setlength{\tabcolsep}{4pt}
\begin{tabular}{@{}clllcc@{}}
\toprule
$k$ & Selection & Blocks kept & Test MASE increase & Better / worse & Speedup \\
\midrule
4 & Prefix & 1--4 & $+3.1\%$ ($+37.0\%$) & & $2.36\times$ \\
 & Uniform & 1, 4, 8, 12 & $+8.9\%$ ($+42.2\%$) & 0 / 10 & $2.35\times$ \\
 & Early + late & 1, 2, 11, 12 & $+10.1\%$ ($+31.3\%$) & 0 / 11 & $2.37\times$ \\
 & Middle & 4--7 & $+7.2\%$ ($+63.3\%$) & 0 / 12 & $2.24\times$ \\
\midrule
6 & Prefix & 1--6 & $+2.3\%$ ($+20.1\%$) & & $1.72\times$ \\
 & Uniform & 1, 3, 5, 8, 10, 12 & $+4.2\%$ ($+17.5\%$) & 4 / 4 & $1.76\times$ \\
 & Early + late & 1--3, 10--12 & $+7.6\%$ ($+23.1\%$) & 1 / 7 & $1.75\times$ \\
 & Middle & 4--9 & $+7.6\%$ ($+52.0\%$) & 1 / 12 & $1.81\times$ \\
\midrule
9 & Prefix & 1--9 & $+2.2\%$ ($+17.1\%$) & & $1.27\times$ \\
 & Uniform & 1, 2, 4, 5, 7, 8, 10--12 & $+0.9\%$ ($+1.1\%$) & 7 / 1 & $1.28\times$ \\
 & Early + late & 1--5, 9--12 & $+0.9\%$ ($+2.1\%$) & 6 / 1 & $1.25\times$ \\
 & Middle & 2--10 & $+3.4\%$ ($+20.2\%$) & 1 / 6 & $1.24\times$ \\
\bottomrule
\end{tabular}
\end{table}

\begin{figure}[t]
\centering
\includegraphics[width=0.75\textwidth]{figures/single_skip_sensitivity.pdf}
\caption{\textbf{Removing a single Chronos-2 block without an adapter.} Top: median test MASE increase over the full model across 14 datasets, with its interquartile range, when one block is skipped and the other 11 run. Bottom: per dataset.}
\label{fig:app_single_skip}
\end{figure}


% Appendix B.4.5: results on model-dataset pairs documented as held out from pretraining (review period).
% Numbers: experiments/make_rebuttal_numbers.py, section C.

\subsubsection{Pairs held out from pretraining}
\label{sec:app_heldout_pairs}

Nine of the 42 model-dataset pairs are documented as held out from pretraining (Table~\ref{tab:provenance}): Chronos-2 on LOOP Seattle, SZ Taxi, M5, and Monash Traffic; TimesFM-3 on LOOP Seattle, SZ Taxi, and M5; and TiRex on M5 and Monash Traffic. Pairs that are not listed are not counted here, because their documentation does not establish that the model never saw the data. Table~\ref{tab:heldout_pairs} repeats the main analyses on the held-out pairs.

The tunnel begins before the final block in 8 of the 9 pairs; the exception is TiRex on Monash Traffic. At the validation-selected entrance, the test loss differs from the final-block readout by a median $+1.3\%$ (range $-8.3\%$ to $+6.9\%$), the same median as over all 42 pairs. Including the not-listed pairs, 22 of 23 pairs enter before the final block. The hard cut fails and alignment recovers most of its gap for all three models. On SZ Taxi and M5, however, the full models are close to the best constant forecast (validation loss 0.89 to 0.91 of the constant forecast's), so small degradations there carry little information, whereas Monash Traffic, where the full models have the most skill (0.35 for Chronos-2 and 0.40 for TiRex), is the held-out dataset that tolerates truncation worst. With 4, 3, and 2 datasets per model, these results agree with the main conclusions but are too few to establish them for unseen data in general.

\begin{table}[t]
\centering
\caption{\textbf{Pairs documented as held out from pretraining.} Tunnel: pairs whose sustained 5\% tunnel begins before the final block. $\Delta$ test: median change in $Q=9$ test loss at the validation-selected entrance relative to the final-block readout. Hard cut, aligned pathway, and readout: median test MASE increase over the full model at 25\%, 50\%, and 75\% of depth on the held-out datasets; all 14 datasets in parentheses for the aligned pathway. TimesFM-3 uses the zero-anchored adapter (offline; Appendix~\ref{sec:app_timesfm3_ridge}).}
\vspace{4pt}
\label{tab:heldout_pairs}
\small
\setlength{\tabcolsep}{3pt}
\resizebox{\textwidth}{!}{%
\begin{tabular}{@{}lcccccc@{}}
\toprule
Model (held-out datasets) & Tunnel & $\Delta$ test & Hard cut & Aligned pathway & (all 14) & Readout \\
\midrule
Chronos-2 (4) & 4/4 & $+1.8\%$ & $+9.3$, $+10.6$, $+6.3\%$ & $+0.9$, $+1.1$, $+1.1\%$ & ($+4.4$, $+2.3$, $+2.2\%$) & $+7.2$, $+5.7$, $+6.0\%$ \\
TimesFM-3 (3) & 3/3 & $+1.3\%$ & $34$, $89$, $100\times$ & $+3.9$, $+4.5$, $+4.4\%$ & ($+10.6$, $+11.3$, $+8.5\%$) & $+9.2$, $+9.6$, $+9.6\%$ \\
TiRex (2) & 1/2 & $+1.9\%$ & $+160$, $+138$, $+132\%$ & $+27.2$, $+15.6$, $+8.3\%$ & ($+13.4$, $+10.6$, $+6.2\%$) & $+40.8$, $+39.4$, $+30.5\%$ \\
\bottomrule
\end{tabular}%
}
\end{table}


\end{document}
